% concatenated from QV_sinebetacorrelation.tex
\documentclass[12pt]{article}
\title{On the pair correlation function of the $\Sineb$ process}
\usepackage{palatino}
\date{}
\author{Yahui Qu and Benedek Valk\'o}


\oddsidemargin 0in \topmargin 0in \headheight 0in \headsep 0in
\textheight 9in \textwidth 6.5in
\renewcommand{\baselinestretch}{1.25}
%\usepackage{units}
\usepackage{amsmath, amsthm, amssymb}
\usepackage{graphicx}
\usepackage{amsmath, enumerate}
\usepackage{color, url}
%\usepackage[notref,notcite]{showkeys}
%\usepackage[longnamesfirst]{natbib}
%\bibpunct[ ]{(}{)}{,}{a}{}{,}

%\bibliographystyle{dcu}
\usepackage{hyperref}
%\RequirePackage[colorlinks,urlcolor=blue,linkcolor=black,citecolor=blue]{hyperref}
%\theoremstyle{theorem}
    \newtheorem{theorem}{Theorem}
    \newtheorem{lemma}[theorem]{Lemma}
    \newtheorem{proposition}[theorem]{Proposition}
    \newtheorem{corollary}[theorem]{Corollary}
    \newtheorem{claim}[theorem]{Claim}
    \newtheorem{question}[theorem]{Question}
    \newtheorem{fact}[theorem]{Fact}
        \newtheorem{conjecture}[theorem]{Conjecture}

\theoremstyle{definition} % For roman text in the body
    \newtheorem{definition}[theorem]{Definition}
    \newtheorem{remark}[theorem]{Remark}
    \newtheorem{example}[theorem]{Example}
    \newtheorem{assumption}{Assumption}
    \newtheorem{algorithm}[theorem]{Algorithm}
    \newtheorem{problem}{Problem}


    
\newcommand{\eps}{\varepsilon}
\newcommand{\Z}{{\mathbb Z}}
\newcommand{\ZZ}{{\mathbb Z}}
\newcommand{\FF}{{\mathcal
F}}
\newcommand{\UU}{{\mathbb U}}
\newcommand{\trb}{{\scriptstyle \frac{2}{\sqrt{\beta}}}}
\newcommand{\R}{{\mathbb R}}
\newcommand{\CC}{{\mathbb C}}
\newcommand{\ud}{{\mathbb U}}
\newcommand{\Rnn}{{\R_{\geq 0}}}
\newcommand{\N}{{\mathbb N}}
\newcommand{\HH}{{\mathbb H} }
\newcommand{\ev}{{\rm   E}}
\newcommand{\pr}{\mbox{\rm P}}
\newcommand{\lstar}{{\raise-0.15ex\hbox{$\scriptstyle \ast$}}}
\newcommand{\ldot}{.}
\newcommand{\vfi}{\varphi}
\theoremstyle{remark} % For an italic header, more subtle than definition style
\newcommand{\Sineb}{\operatorname{Sine}_{\beta}}
\newcommand{\Sine}{\operatorname{Sine}}
\newcommand{\Airyb}{\operatorname{Airy}_{\beta}}
\newcommand{\Bessel}{\operatorname{Bessel}}
\newcommand{\Sch}{\operatorname{Sch}}
\newcommand{\HPb}{\operatorname{HP}_{\beta, \delta}}
\newcommand{\HPbb}{\operatorname{HP}_{\beta, \beta/2}}
\newcommand{\Sineop}{\mathtt{Sine}_{\beta}}
\newcommand{\Circop}{\mathtt{Circ}_{\beta,n}}
\newcommand{\Besselop}{\mathtt{Bessel}_{\beta,a}}
\newcommand{\Airyop}{\mathtt{Airy}_{\beta}}
\newcommand{\Huaop}{\mathtt{HP}_{\beta,\delta}}
\newcommand{\Schop}{\mathtt{Sch}_{\nu}}
\newcommand{\Dirop}{\mathtt{Dir}}
\newcommand{\MT}{\mathrm{T}}
\newcommand{\MR}{\mathrm{R}}
\newcommand{\MS}{\mathrm{S}}
\newcommand{\half}{\nicefrac{1}{2}\,}
\newcommand{\gl}{\lambda}
\newcommand{\Dfi}{\Delta \vfi}
\renewcommand{\O}{\mathcal{O}}
\definecolor{violet}{rgb}{0.8,0,0.2}
\newcommand{\cd}{\stackrel{d}{\longrightarrow}}
\newcommand{\ed}{\stackrel{d}{=}}
% \newcommand{\cd}{\stackrel{d}{\Rightarrow}}
\newcommand{\cP}{\stackrel{P}{\to}}
\newcommand{\cB}{{\mathcal B}}
\newcommand{\cW}{{\mathcal W}}
\newcommand{\cV}{{\mathcal V}}

\newcommand{\cA}{{\mathcal A}}

\newcommand{\cE}{{\mathcal E}}
\newcommand{\cH}{{\mathcal H}}
\newcommand{\cL}{{\mathcal L}}
\newcommand{\cG}{{\mathcal G}}
\newcommand{\btau}{{\boldsymbol{\tau}}}
\newcommand{\T}{\dagger}

\newcommand{\ttHP}{\mathtt{HP}}

\newcommand{\Tr}{{\operatorname{Tr\,}}}
\newcommand{\mat}[4]{\left[ \hspace{-0.4em} \begin{array}{cc}
#1 & #2  \\
#3 & #4  \\
\end{array} \hspace{-0.4em}\right]}
\newcommand{\bin}[2]{\binom{#1}{#2}}
% \newcommand{\bin}[2]{\left[
% \begin{array} {c}
% #1 \\
% #2
% \end{array}
% \right]}
\newcommand{\la}{\langle}
\newcommand{\ra}{\rangle}
\newcommand{\SU}{\operatorname{SU}(1,1)}
\newcommand{\SL}{\operatorname{SL}_2(\mathbb R)}
\newcommand{\su}{\mathfrak{su}(1,1)}
\newcommand{\ac}{{\text{\sc ac}}}
\newcommand{\dom}{\operatorname{dom}}
\newcommand{\range}{\operatorname{Range}}
\newcommand{\tr}{\operatorname{tr}}
\newcommand{\ind}{\mathbf 1}
\newcommand{\tch}{\upsilon} % the time change function

\newcommand{\pochup}[2]{(#1)^{\overline{#2}}}

\newcommand{\pochdown}[2]{(#1)^{\underline{#2}}}

\newcommand{\coeff}{c}
\newcommand{\sinc}{\operatorname{sinc}}

\newcommand{\CbE}{\mathtt{C}\beta\mathtt{E}}

\newcommand{\CJbE}{\mathtt{CJ}\beta\mathtt{E}}

\definecolor{darkgreen}{rgb}{.2,0.4,0.3}

\newcommand{\benedek}[1]{\textcolor{red}{\tt{#1}}}


\newcommand{\yahui}[1]{\textcolor{blue}{\tt{#1}}}

\newcommand{\todo}[1]{\textcolor{darkgreen}{\tt{To do: #1}}}


\newcommand{\szeta}{\boldsymbol{\zeta}}

\newcommand{\tph}{\tilde \phi}

\newcommand{\mtr}[1]{\mathrm{#1}}

\newcommand{\tl}{\tilde}
\newcommand{\res}{{\mathtt{r}\,}}

\newcommand{\ress}{{\mathtt{r}}}


\newcommand{\q}{q}

\newcommand{\qq}{\mathbf{q}}

\newcommand{\spec}{\operatorname{spec}}


\newcommand{\sech}{\operatorname{sech}}

%\usepackage{natbib}
\bibliographystyle{plain}

\begin{document}
\maketitle

\begin{abstract}
We study the $\Sineb$ process, the bulk point process scaling limit of beta-ensembles. We provide a representation of  its pair  correlation function for all $\beta>0$ via a stochastic differential equation. 
We show that the pair correlation function is continuous in $\beta$, and provide estimates for its asymptotic decay.  We
    recover the classical explicit formula for the pair correlation function in the $\beta=2$ and $4$ cases. For  $\beta=2n$, we derive the power series expansion of the pair correlation function, and express it 
    in terms of a size $n$ linear ordinary differential equation system. 
    We obtain our results by studying the density of the  $\HPb$ process, the point process limit of the circular Jacobi beta-ensembles. %, which was derived in \cite{LV_HP} using the stochastic operator framework. 
\end{abstract}


\section{Introduction}

% For a given $\beta>0$, the size $n$ circular beta-ensemble ($\CbE_n$) is the distribution of $n$ points $e^{i \theta_j}, 1\le j\le n, \theta_j\in (-\pi,\pi]$ on the unit circle given by the following joint density function
% \begin{align}\label{eq:cbe}
%     \frac{1}{Z_{\beta, n}^{C}} \prod_{1\le j<k\le n} |e^{i \theta_j}-e^{i \theta_k}|^\beta.
% \end{align}
% Here $Z_{\beta, n}^{C}$ is an explicitly computable constant. We say that $\Lambda_n\sim \CbE_n$ if XXX. 
% For $\beta=2$ this distribution is the same as the joint distribution of the eigenvalues of a uniformly chosen $n\times n$ unitary matrix. 

The $\Sineb$ process is the bulk scaling limit of beta-ensembles. It is a translation invariant process with density $\tfrac1{2\pi}$. It was first introduced as the bulk scaling limit of the Gaussian beta-ensemble for general $\beta>0$ in \cite{BVBV}. Later it was shown to be the same as the point process limit of the circular beta-ensemble (\cite{BVBV_sbo},\cite{Nakano}) which was first derived by Killip and Stoiciu in \cite{KS}. Finally, \cite{BEY} proved that $\Sineb$ is the universal bulk scaling limit of a general class of beta-ensembles. 


For a simple point process $\Xi$ on $\R$ and a Borel set $A\subset \R$, we denote by $\Xi(A)$ the number of points of $\Xi$ in $A$. 
The $k$-point correlation (or joint intensity) function of $\Xi$ (if it exists) is the symmetric function $\rho^{(k)}:\R^k\to \R_+$ that satisfies the following: if  $D_1, \dots, D_k\subset \R$ are disjoint Borel sets then
% When the $k$-point correlation (or joint intensity) function of $\Xi$  exists, it is the symmetric function $\rho^{(k)}:\R^k\to \R_+$ that satisfies the following: if  $D_1, \dots, D_k\in \R$ are disjoint Borel sets then
\begin{align}
E[\prod_{j=1}^k \Xi(D_j)]=\int_{D_1} \cdots  \int_{D_k} {\rho^{(k)}(x_1,\dots,x_k)} dx_k\, \dots dx_1.
\end{align}
The function $\rho^{(1)}$ (if it exists) is just the Radon–Nikodym derivative of the expected counting function $\nu(A)=E[\Xi(A)]$ with respect to the Lebesgue measure, this is the density of the point process $\Xi$. 


Many of the classical results of random matrix theory use $k$-point correlation functions to identify the point process limit of various random matrix ensembles. For the classical values $\beta=1,2,4$, the finite beta-ensembles possess special algebraic structures. These processes are determinantal (for $\beta=2$) or Pfaffian (for $\beta=1, 4$), which means that it is possible to express their correlation functions in terms of determinantal or Pfaffian formulas built from the orthogonal polynomials of the reference measure of the ensemble. For circular and Gaussian beta-ensembles, one can take limits of these formulas to obtain explicit expressions for the $k$-point correlation of the $\Sineb$ process for $\beta=1,2,4$  (see \cite{AGZ, mehta, ForBook}). 


% When $\beta=2n$ (i.e.~$\beta$ is an even integer) Forrester 
%  provides a representation for the correlation functions of the circular beta-ensemble in terms of generalized hypergeometric functions and Jack polynomials (see Chapter 13 of \cite{ForBook}), and evaluates the scaling limits of these functions. 


In contrast, the $\Sineb$ process           for general $\beta>0$ is  characterized via its counting function using a one-parameter family of stochastic differential equations (\cite{BVBV}, \cite{KS}), or as the spectrum of a random differential operator (\cite{BVBV_sbo}). It has not been clear how one can obtain the classical $\beta=1,2,4$ results from these characterizations of $\Sineb$, and whether it is possible to get any information about the correlation functions. Our paper provides the first positive results in this direction, by studying the pair (or two-point) correlation function $\rho^{(2)}_\beta(x,y)$ of the $\Sineb$ process. 


\subsection*{Results for general $\beta>0$}

Our first result provides a characterization of $\rho^{(2)}_\beta(0,\lambda)$ in terms of a stochastic differential equation. Since the $\Sineb$ process  is translation and reflection invariant, we have $\rho^{(2)}_\beta(x,y)=\rho^{(2)}_\beta(0,|y-x|)$. 
Here, and in the rest of the paper, for $x\in \R$ and a non-negative integer $k$ we use the following notation for the rising factorial (or Pochhammer symbol):
\begin{align}\label{eq:Poch}
    (x)^{\uparrow k}=\prod_{j=0}^{k-1}(x+j), 
\end{align}
with the empty product defined as 1. 
\begin{theorem}\label{thm:sineb_corr}
    Let $\beta>0$, $\lambda>0$, and consider the unique strong solution $\alpha_{\lambda}(u)$ of the following stochastic differential equation on $(-\infty,0]$:
    \begin{align}
    \label{eq:alphaSDE_b}
	d\alpha_\lambda = \lambda\tfrac{\beta}{4}e^{\frac{\beta}{4}u} du -\frac{\beta}{2} \sin(\alpha_\lambda(u)) du+ \Re[(e^{-i\alpha_\lambda(u)}-1)dZ], \qquad \lim_{u\to -\infty} \alpha_\lambda(u)=0. 
\end{align}
    The  pair correlation function of the $\Sineb$ process is given by
\begin{align}\label{eq:twopoint_cor}
        \rho_\beta^{(2)}(0,\lambda)=\frac{1}{4\pi^2}+\frac{1}{2\pi^2}\sum_{k=1}^\infty  \frac{(-\beta/2)^{\uparrow k}}{(1+\beta/2)^{\uparrow k}} E[\cos(k \alpha_\lambda(0)].
    \end{align}  
\end{theorem}
For $\delta>0$ one has the uniform bound
\begin{align}
\left|\frac{(-\delta)^{\uparrow k}}{(1+ \delta)^{\uparrow k}}\right|\le c k^{-1-2  \delta},
   \label{eq:coeff_bnd} 
\end{align}
with a $\delta$-dependent  constant $c$, hence \eqref{eq:twopoint_cor} gives an absolutely convergent sum. 

One can realize the solutions of \eqref{eq:alphaSDE_b} for all $\beta>0$ and $\lambda>0$ on the same probability space, in a way that the solution is almost surely continuous in its parameters. This allows us to prove the following regularity result. 
\begin{theorem}\label{thm:beta_cont}
    The function $\rho_{\beta}^{(2)}(0,\lambda)$ is continuous in both $\beta$ and $\lambda$.  For $\beta>2, \lambda>0$ the function is jointly continuous in $\lambda$ and $\beta$. 
\end{theorem}
Identifying the large $\lambda$ behavior of 
the  pair correlation function $\rho^{(2)}_\beta(0,\lambda)$ for general $\beta>0$ has been a long-standing open problem. Using the representation of Theorem \ref{thm:sineb_corr} we are able to provide the following bound on the asymptotic decay of the truncated correlation function $\rho^{(2)}_\beta(0,\lambda)-\rho^{(1)}_\beta(\lambda)^2=\rho^{(2)}_\beta(0,\lambda)-\frac{1}{4\pi^2}$. 
\begin{theorem}[Decay of the truncated pair correlation]\label{thm:large_lambda}
There is a $\beta$-dependent constant $c$ so that we have the following bound for $\lambda\ge 2$: 
\begin{align}\label{eq:large_lambda}
     \left|\rho^{(2)}_\beta(0,\lambda)-\frac{1}{4\pi^2}\right|\le c\left(\lambda^{-1}+\lambda^{-\beta/2}+\lambda^{-4/\beta}\right).
\end{align} 
\end{theorem}
For $\beta\ge 4$, the main term in our upper bound is $\lambda^{-4/\beta}$, which matches the conjectured order of the truncated pair correlation function in this regime (see equation \eqref{eq:rho2_2n_asympt} below, and the discussion around it). While preparing this manuscript, we have learned from Laure Dumaz and Martin Malvy \cite{DM2025} that they have also obtained polynomially decaying bounds on the truncated pair correlation function of $\Sineb$, with methods that are different from ours. 

\subsection*{Results for $\beta=2n$}

The beta-ensembles do not possess determinantal/Pfaffian structure for values other than the classical values 1, 2, and 4. However,  Forrester derived exact formulas  the circular beta-ensemble for $\beta=2n$, $n\in \ZZ_+$ (see \cite{Forrester_1992}, \cite{Forrester_1994}, and Chapter 13 of \cite{ForBook}). These formulas lead to an $n$-variable integral representation of the pair correlation function of the $\Sineb$ process when $\beta=2n$. We will review these results in more detail in the \textit{Historical background} section below.

We also present some new results for $\beta=2n$.  
The expression $(-x)^{\uparrow k}$ is zero if $x$ is a positive integer and $k>x$. This means that the expression \eqref{eq:twopoint_cor} in Theorem \ref{thm:sineb_corr} becomes a finite sum if $\beta=2n$. In this case we 
can provide more explicit representations of the pair correlation function.  We show that  $\rho^{(2)}_{2n}(0,\lambda)$ is analytic in $\lambda^2$,  express it in terms of the solution of an $n$-dimensional linear ODE system, and also as a power series with coefficients obtained from a matrix-valued recursion. For the full result see Theorem \ref{thm:even_beta}, Corollary \ref{cor:int_case_rho}, and the discussion following these statements. Here, we only present in detail the result of the power series representation.

Fix $n\in \Z_+$.  
Let $\mathbf{A}_{n}$ be a tridiagonal  and $\mathbf{B}_n$ a diagonal $n\times n$ matrix with nonzero entries given by 
\begin{align}\label{eq:AB}
      [\mathbf{A}_n]_{k,k}=-k^2, \quad [\mathbf{A}_n]_{k,k-1}=\frac12 k(k+n), \quad [\mathbf{A}_n]_{k,k+1}=\frac12 k(k-n),\quad [\mathbf{B}_n]_k=k.
\end{align}
Let $\mathbf{e}_n, \mathbf{v}_n\in \R^n$ be defined as 
\begin{align}\label{eq:vectors}\mathbf{e}_n=[1,0,\dots,0]^T, \qquad 
[\mathbf{v}_n]_k=(-1)^k\frac{ \binom{2n}{n+k}}{\binom{2n}{n}}, \qquad 1\le k\le n.
\end{align}

\begin{theorem}[Power series representation for $\rho_{2n}^{(2)}(0,\lambda)$]\label{thm:sineb_series}
Define $\mathbf{s}_k \in \CC^n$ recursively as
\begin{align}\label{eq:q_coeff}
    \mathbf{s}_0&=-\frac{n+1}{2}\mathbf{A}_n^{-1}\mathbf{e}_n, \qquad \mathbf{s}_k=i \left(k \mathbf{I}-\tfrac{2}{n} \mathbf{A}_n\right)^{-1} \mathbf{B}_n \mathbf{s}_{k-1}, \qquad k\ge 1.
\end{align}
Then 
\begin{align}\label{eq:rho2_2n}
    \rho_{2n}^{(2)}(0,\lambda)=\frac{1}{2\pi^2} \sum_{j=1}^\infty \mathbf{v}_n^{T} \cdot \mathbf{s}_{2j}\,  \lambda^{2j}.
\end{align}    
\end{theorem}
As part of the proof, we will show that the sequence $\mathbf{s}_k, k\ge 0$ is well-defined, and the power series converges for all $\lambda\in \R$. 

Theorem \ref{thm:sineb_series} allows us to recover the well-known Sine-kernel for the $\beta=2$ case. Indeed, if  $\beta=2$ (i.e.~$n=1$), then we have $\mathbf{A}_1=-1$, $\mathbf{B}_1=1$, $\mathbf{v}_1=-\tfrac{1}{2}$,  $\mathbf{s}_k=\tfrac2{(k+2)!} i^k $, hence \eqref{eq:rho2_2n} yields
\begin{align}\label{eq:sine_kernel}
    \rho_{2}^{(2)}(0,\lambda)=\frac{1}{4\pi^2}\left(1-\frac{\sin^2(\tfrac{\lambda}{2})}{(\lambda/2)^2}\right).
\end{align}
In Section \ref{sec:additional} we will find a closed formula for the function $\sum_{k=0}^\infty \mathbf{s}_{j}  \lambda^{j}$ in the $\beta=4$ case, which allows us to recover the known expression (see \eqref{eq:beta=4} below) for the pair correlation function of the $\Sine_4$ process.

As an additional application, we recover the asymptotics of the pair correlation function near 0 in the $\beta=2n$ case. (This also follows directly from the integral representation of Forrester, see the discussion below.) 

\begin{corollary}\label{cor:small_lambda}
  We have the following asymptotics for $\lambda$ near 0:  
  \begin{align}\label{eq:small_lambda}
      \rho_{2n}^{(2)}(0,\lambda)= \frac1{4\pi^2} \cdot \frac{n^{2n} (n!)^3 }{(2n!)(3n!) }\lambda^{2n}+\mathcal{O}(\lambda^{2n+2}).
  \end{align}
  
\end{corollary}



\subsection*{Strategy of the proof}

Our starting point is the following observation from \cite{BVBV_Palm}: if one can properly define the $\Sineb$ process conditioned to have a point at 0, then the one-point function of the conditioned process is  $\frac{\rho_\beta^{(2)}(0,x)}{\rho_\beta^{(1)}(x)}$. 

The $\Sineb$ process is the point process limit of the circular beta-ensemble. We can readily condition a size  $n$ circular beta-ensemble to have a point at angle 0. After removing this point, we obtain the size $n-1$ circular Jacobi beta-ensemble with parameter $\delta=\beta/2$. The scaling limit of the circular Jacobi beta-ensemble was obtained in \cite{LV_HP}, this is the $\HPb$ process. (See Section \ref{sec:background} for precise statements.) We will show (using the results of \cite{BVBV_Palm}) that the $\HPb$ process with $\delta=\beta/2$ is the same as the $\Sineb$ process conditioned to have a point at 0, with that point removed (see Proposition \ref{prop:Palm} below). In particular, with $\rho_{\beta,\delta}^{(1)}(\lambda)$ denoting the density of $\HPb$, we will show that  
\begin{align}\label{eq:palm}
  \rho_\beta^{(2)}(0,\lambda)=  \rho_{\beta,\beta/2}^{(1)}(\lambda){\rho_\beta^{(1)}(\lambda)}=\frac{1}{2\pi} \rho_{\beta,\beta/2}^{(1)}(\lambda).
\end{align}
\cite{LV_HP} provides a characterization of $\HPb$   as the result of a shooting problem involving an SDE similar to \eqref{eq:alphaSDE_b}. Using this characterization, we can prove a version of Theorem \ref{thm:sineb_corr} for $\rho_{\beta,\delta}^{(1)}(\lambda)$ with $\beta>0, \delta>0$, see Theorem \ref{thm:1p_HP} below. Theorem \ref{thm:sineb_series} and Corollary \ref{cor:small_lambda} will also follow from corresponding results on the density of the $\HPb$ process.


\subsection*{Historical background}

As mentioned before, for $\beta=2$, the circular and the circular Jacobi beta-ensembles are determinantal, the $k$-point correlation functions can be obtained as the determinant of a $k\times k$ matrix built from a kernel function. Taking the scaling limit of these formulas, one obtains exact determinantal expressions for all correlation functions of $\Sineb$ and $\HPb$ for $\beta=2$, in particular, one obtains \eqref{eq:sine_kernel}. For $\beta=1$ and $\beta=4$, one can write down the correlation functions of the finite and the limiting ensembles in terms of Pfaffian formulas (\cite{mehta}, \cite{AGZ}, \cite{ForBook}). For $\beta=4,$ this leads to the following expression for the pair correlation function of the $\Sineb$ process:
\begin{align}\label{eq:beta=4}
    % &\rho_{\beta=2}^{(2)}(0,\lambda) = \frac{1}{4\pi^2}-\frac{1}{2\pi^2}\frac{\sin(\tfrac{x-y}{2})}{x-y};\\
     &\rho_{4}^{(2)}(0,\lambda) = \frac{1}{4\pi^2}\left(1- \sinc^2(\lambda) + \sinc'(\lambda)\int_0^{\lambda} \sinc(t)dt\right),\qquad \sinc(x)=\frac{\sin(x)}{x}.
     % &\rho_{\beta=4}^{(2)}(0,\lambda) = \frac{1}{4\pi^2}\left(1- \frac{1}{4\pi^2} \left(\frac{\sin(x-y)}{x-y}\right)^2 + \frac{1}{4\pi^2}\frac{\partial}{\partial x} \frac{\sin(x-y)}{x-y}\int_0^{x-y} \frac{\sin t}{t}dt\right)
\end{align}
% (7.2 and 7.92)
% \benedek{write it in terms of $\operatorname{sinc}(x):=\frac{\sin(x)}{x}$}
 When $\beta=2n$ (i.e.~$\beta$ is an even integer), Forrester 
  provides a representation for the correlation functions of the circular beta-ensemble in terms of generalized hypergeometric functions and Jack polynomials (see \cite{Forrester_1992}, \cite{Forrester_1994}, and Chapter 13 of \cite{ForBook}), and evaluates the scaling limit of these functions. 
% For $\beta=2n, n\in \Z_+$, Chapter 13 of \cite{ForBook} provides an expression for the bulk scaling limit of the correlation functions of the circular beta-ensemble in terms of generalized hypergeometric functions. 
We record the result for the limit of the pair correlation function here, which, if one can justify the exchange of limits,  gives $\rho_{2n}^{(2)}(0,\lambda)$:
\begin{align}\notag
 &\frac1{4\pi^2} \cdot \frac{n^{2n} (n!)^3 }{(2n!)(3n!) } \frac{e^{-i n \lambda} \lambda^{2n}}{S_{2n}(-1+1/n,-1+1/n, 1/n)}\times\\
 &\qquad \int_{[0,1]^{2n}} \prod_{j=1}^{2n} \left(e^{i \lambda u_j} u_j^{-1+1/n} (1-u_j)^{-1+1/n} \right)\prod_{j<k}|u_j-u_k|^{2/n} \prod_{j=1}^{2n}du_j .\label{eq:paircorr_For}
\end{align}
Here $S_{2n}(-1+1/n,-1+1/n, 1/n)$ is the Selberg integral on $[0,1]^{2n}$ that one obtains by setting $\lambda=0$ in the integral in the second line of \eqref{eq:paircorr_For}. Note we can immediately identify the constant in the asymptotics \eqref{eq:small_lambda} by taking $\lambda\to 0$. It would be interesting to see if one could recover the power series representation \eqref{eq:rho2_2n} directly from this result, or vice versa.  


% should be possible to recover the power series representation of Theorem  \ref{thm:sineb_series} directly from \eqref{eq:paircorr_For}, \benedek{see open problems}


By analyzing the integral in \eqref{eq:paircorr_For},
Forrester provides the following large $\lambda$ asymptotics of \eqref{eq:paircorr_For} for $\beta=2n$:
\begin{align}\notag
4\pi^2    \rho_{\beta}^{(2)}(0,\lambda)=&1-\frac{4}{\beta \lambda^2}+\mathcal{O}(\lambda^{-4})\\
&+\sum_{k=1}^{n}\frac{a_k}{\lambda^{4k^2/\beta}}\left(\cos(k \lambda)+b_k\frac{\sin(k\lambda)}{\lambda}+\mathcal{O}\left(\frac{|\cos(k\lambda)|}{\lambda^2}\right)\right)\label{eq:rho2_2n_asympt},
\end{align}
with explicitly given coefficients $a_k, b_k$ (see Proposition 13.13, \cite{ForBook}).
Similar asymptotics, with an infinite sum of oscillating terms in \eqref{eq:rho2_2n_asympt}, 
were obtained by Haldane \cite{Haldane1981}
in the study of one-component, one-dimensional quantum fluids. It is conjectured that Haldane's asymptotics hold 
for the pair correlation $\rho_{\beta}^{(2)}(0,\lambda)$ for general $\beta>0$ as well (see e.g.~\cite{Forrester_1984}). Note that the largest order term in this conjectured expansion is of order $\lambda^{-\min(2,4/\beta)}$. Our bound in Theorem \ref{thm:large_lambda} matches this order for $\beta\ge 4$. 


In \cite{Forrester_20212} (building on the results of \cite{FR2012}) Forrester outlines how one can represent $\rho_{2n}^{(2)}(0,\lambda)$ as the solution of an order $2n+1$ differential equation, and provides these representations for $\beta=1,2,4,6$. \cite{Forrester_2021} gives similar results for the scaling limit of the one-point function of the circular Jacobi beta-ensemble.  






% Following Haldane's 

% ground state pair correlations of general one-component, one-dimensional quantum fluids

% \benedek{Add the Haldane conjecture discussion here, and make the connection about the order of decay in \eqref{eq:large_lambda} } 

% \yahui{I found this paper:Coulomb and Riesz gases: The known and the unknown Page 164 very helpful but i dont knwo if i can just say what the author said, and i tried to read the paper he quoted they r all abt physics which is kinda hard for me to understand.} 


% When $\beta$ is a rational number \cite{okounkov1997} offers a formula for the $k$-point correlation function of the circular beta ensemble using the Selberg integral, see also \cite{Forrester_1995} where a conjectured formula is presented for the pair correlation function.


 




% From our random operator approach, we can recover the formula in \cite[Chapter~7]{ForBook}for $\beta=2,4$, and 
% we give exact solvable formulas for the pair correlation function in the case of even $\beta>0$. For general $\beta>0$, we can give a generic formula (\ref{eq:twopoint_cor}) expressed using the $\alpha_\lambda(u)$ defined as above \ref{eq:alphaSDE}, with which we will be able to prove the results 




The random operator approach to study the point process limits of beta-ensembles was initiated by Dumitriu-Edelman \cite{DE} and Edelman-Sutton \cite{ES}. Building on their work, the point process limits of beta-ensembles have been characterized as the spectrum of random differential operators, or by describing the counting function using systems of stochastic differential equations \cite{RRV}, \cite{RR}, \cite{KS}, \cite{BVBV}, \cite{BVBV_sbo}, \cite{LV_HP}. These descriptions provided tools to study various asymptotic quantities related to these processes (e.g.~large gap asymptotics \cite{BVBV2}, \cite{RRZ}, asymptotics for the counting function \cite{KVV}, \cite{HP2018}, large deviations for the number of points in a small or a large interval \cite{HV}, \cite{HV2}), but before the current paper there have been only a few results that recovered the exact formulas known for the classical beta values. The most significant result where this is actually achieved  is from \cite{Spike1}, where  the authors recovered explicit formulas for $\beta=2,4$ for the soft-edge limit of the Gaussian beta-ensemble, providing an independent proof of the Painlev\'e representations of 
Tracy-Widom distributions originally given in \cite{TW1994}, \cite{TW1996}.

There are several other promising approaches to study point process limits of beta-ensembles for general $\beta>0$. \cite{DLR2021} gives a characterization of the $\Sineb$ process using the Dobrushin-Lanford-Ruelle formalism, \cite{Leble2021} uses these results to prove a CLT for the linear statistics for $\Sineb$. 
The method of loop equations (introduced in \cite{Johansson1998} provides a powerful tool to study finite beta-ensembles, see e.g.~\cite{Lambert2021} for an application.


We also mention here the unpublished preprint  \cite{okounkov1997}, where the limit of the $k$-point correlation function of the circular beta-ensemble is presented in the case when $\beta$ is a rational number. 





% There have been several results about the scaling limit of the pair correlation for the $\Sineb$ process for special values of $\beta$.\\
% In terms of the exact formulas, when $\beta=1,2,4$, the eigenvalues of Gaussian beta ensemble form either a determinantal($\beta=2$) process or a  Pfaffian($\beta=1,4$) process.(\cite{Dyson62_1},\cite{Tracy_1998}). 

% \benedek{state the formulas here}

% Exploiting the structures of these processes, the exact formulas for the pair correlations are given in \cite{mehta},\cite{ForBook}. When $\beta=2n,\; n\in \mathbf{N}$ is even, Forrester gave an integral expression for the scaling limit of the pair correlation function by studying the Jack polynomials and the Selberg integral. He also provided the large $\lambda$ asymptotic of the correlation function \cite{ForBook}. When $\beta = \frac{p}{q}$, $p, q \in \mathbf{N}$ is rational, Okounkov obtained a formula for the n-point correlation function of the circular beta ensemble using the Selberg integral \cite{okounkov1997npointcorrelationsloggasrational}.\\ 
% There is also other representations of the correlation function. In \cite{Forrester_20212}, the author showed for $\beta=1,2,4,6$, the 2- point correlation function for Circular beta ensemble satisfies a differential equation of order $\beta+1$. In \cite{Forrester_2021}, the authors showed the density function of circular Jacobi beta ensemble satisfies certain differential equation for $\beta=1,2,4$. Later, we will see that scaling limit of the density function of the circular Jacobi beta ensemble corresponds to the scaling limit 2-point correlation function for circular beta ensemble. \\

% \yahui{Conjecture for decay I need to know more abt those physics papers}


% \benedek{State the conjecture for the decay}


\subsection*{Outline of the paper}

In Section \ref{sec:background}, we review some known results on the circular and circular Jacobi beta-ensembles and their point process limits. We also sketch the proof of equation \eqref{eq:palm}. Section \ref{sec:1point} contains the proof of Theorem \ref{thm:sineb_corr} and the analogous result on the density of the $\HPb$ process with $\delta>0$.  Section \ref{sec:large_lambda} provides the proof for the large $\lambda$ asymptotics of $\rho_{\beta}^{(2)}(0,\lambda)$, and Section  \ref{sec:delta_n} provides the results for the $\beta=2n$ case, with corresponding results about $\HPb$ with $\delta=n$. 
Section \ref{sec:additional} provides some additional results, including the proof of Theorem \ref{thm:beta_cont}. 
Finally, Section \ref{sec:open} lists some open problems. 









\subsection*{Acknowledgments} 

The authors  thank Laure Dumaz, Martin Malvy and Yun Li for useful discussions. 
B.V.~was partially supported by  the University of Wisconsin – Madison Office of the Vice Chancellor for Research and Graduate Education with funding from the Wisconsin Alumni Research Foundation and by the National Science Foundation award DMS-2246435. 


\section{Preliminaries on beta-ensembles}\label{sec:background}


\subsection{The circular and the circular Jacobi beta-ensemble and their limits}



Let $\beta>0$, $n\in \Z_+$, and $\delta\in \CC$ with $\Re \delta>-1/2$. The size $n$ circular Jacobi beta-ensemble with parameters $\beta, \delta$ is the joint distribution  of 
$n$  points $\{e^{i \theta_1}, \dots, e^{i \theta_n}\}$ on the unit circle $\{|z|=1\}$, where the joint density function of the angles $\theta_j\in [-\pi,\pi)$, is given by
\begin{align}
 \frac{1}{Z_{n,\beta,\delta}^{\textup{CJ}}} \prod_{1\le j<k\le n} \left|e^{i \theta_j}-e^{i \theta_k}\right|^\beta  \prod_{k=1}^n (1-e^{- i \theta_k})^{ \delta}(1-e^{ i \theta_k})^{ \bar\delta} \label{eq:PDF_cjacobi}, \qquad \theta_j\in [-\pi,\pi).
\end{align}
Here  $Z_{n,\beta,\delta}^{\textup{CJ}}$ is an explicitly computable normalizing constant. For $\delta=0$ the distribution becomes the circular beta-ensemble. We  write $\Lambda_n\sim \textup{CJ}_{n,\beta,\delta}$ to denote that the random set $\Lambda_n=\{\theta_1, \dots, \theta_n\}$ has  joint density given by \eqref{eq:PDF_cjacobi}, and we use the notation $\textup{C}_{n,\beta}$ for $\textup{CJ}_{n,\beta,0}$, the size $n$ circular beta-ensemble.


Killip and Stoiciu \cite{KS} showed that if $\Lambda_n\sim \textup{C}_{n,\beta}$ then $n\Lambda_n$ converges to a point process as $n\to \infty$. The point process limit was characterized via a one-parameter family of stochastic differential equations via a shooting problem.  This process was later shown to be the same as the   $\Sineb$ process, the bulk scaling limit of the Gaussian beta-ensemble \cite{BVBV_sbo}, \cite{Nakano}, which can also be characterized as the spectrum of a random differential operator.

In \cite{LV_HP} it vas shown that if $\Lambda_n\sim \textup{CJ}_{n,\beta,\delta}$ then $n \Lambda_n$ converges to a point process $\HPb$. (Note that $\HPb$ with $\delta=0$ is just $\Sineb$.) This process can also be characterized in multiple ways (as the spectrum of a certain random differential operator, as the zero set of a certain random entire function). We review the characterization that will be useful for our purposes. 

Let $Z=d\mathbf{B}_1+i d\mathbf{B}_2$ be a two-sided complex Brownian motion with standard independent real and imaginary parts. 
For a given $\beta>0, \Re \delta>-1/2$ consider the unique strong solution $\alpha_\lambda(t)$  of the following one-parameter family of stochastic differential equations:  
\begin{align}
    \label{eq:alphaSDE}
	d\alpha_\lambda = \lambda\tfrac{\beta}{4}e^{\frac{\beta}{4}u} du + \Re[(e^{-i\alpha_\lambda(u)}-1)(dZ-i\delta du)],\qquad \lambda\in \R, t\in(-\infty,\infty) 
\end{align}
with the initial condition
\begin{align}\label{eq:SDE_initial}
    \lim_{u\to-\infty} \alpha_\lambda(u)=0, \qquad \lambda\in \R. 
\end{align}
Let $\Theta$ be a random variable on $[0,2\pi)$ with probability density function 
\begin{align}\label{eq:Theta}
  {\frac{1}{2\pi}} \frac{\Gamma(1+\delta)\Gamma(1+\bar\delta)}{\Gamma(1+\delta+\bar\delta)}(1-e^{-i \theta})^\delta  (1-e^{i \theta})^{\bar \delta}, \qquad \theta\in [0,2\pi),
\end{align}
independent of $Z$.
Then \cite{LV_HP} proved that
\begin{align}\label{eq:HP}
    \HPb\ed\{\lambda\in \R: \alpha_\lambda(0)=\Theta \mod 2\pi\}.
\end{align}
For $\delta=0$ this recovers the characterization of $\Sineb$ given in \cite{KS},  in that case $\Theta$ is uniform on $[0,2\pi)$. Note that the diffusion \eqref{eq:alphaSDE_b} is the same as  \eqref{eq:alphaSDE} with initial condition \eqref{eq:SDE_initial} and $\delta=\beta/2$.

Proposition 9 of \cite{BVBV} provides an exponential tail bound on $\Sineb[0,\eps]$ (i.e.~the number of points in $[0,\eps]$) from which it follows that the $k$th moment of $\Sineb[0,\eps]$ is bounded by a constant multiple of $\eps^k$ as $\eps\to 0$. From this it follows that the $k$th moment measure is absolutely continuous with respect to the Lebesgue measure, and this implies the existence of the $k$th correlation  function for the stationary process $\Sineb$ for all $k\ge 1$ (see \cite{DVJ1}).


\subsection{Properties of the $\alpha_\lambda$ diffusion}


The following proposition summarizes some of the useful properties of the diffusion \eqref{eq:alphaSDE},  see \cite{LV_HP} for additional details and proofs.





\begin{proposition}\label{prop:alphaSDE}Fix $\beta>0$ and $\Re \delta>-1/2$. 


\begin{enumerate}
    \item For any $\nu\in \R$ the system \eqref{eq:alphaSDE} has a unique strong solution $\alpha_{\lambda,\nu}(u)$ on $[\nu,\infty)$ with initial condition $\alpha_{\lambda,\nu}(\nu)=0$. For any given $\nu< u$, with probability one the function $\lambda\mapsto \alpha_{\lambda,\nu}(u)$ is analytic and strictly increasing with $\alpha_{0,\nu}(u)=0$ 

% \item Suppose that $\alpha$ and $\tilde \alpha$ are both strong solutions of \eqref{eq:alphaSDE} on $[u_0,\infty)$. Let $0\le \lambda_1\le \lambda_2$. Then if $P(\alpha_{\lambda_1}(u_0)\le \tilde \alpha_{\lambda_2}(u_0))=1$ then  $\alpha_{\lambda_1}(u)\le \tilde \alpha_{\lambda_2}(u)$ for all $u_0\le u$ with probability one. 
% \benedek{Maybe we don't need this?}

    \item There is a unique strong solution $\alpha_\lambda(u)$ of \eqref{eq:alphaSDE} on $(-\infty,\infty)$ satisfying the initial condition \eqref{eq:SDE_initial}, which  can be obtained as the a.s.~limit $\alpha_\lambda(u)=\lim\limits_{\nu\to -\infty} \alpha_{\lambda,\nu}(u)$. The limit is monotone increasing for $\lambda>0$ and monotone decreasing for $\lambda<0$. 
    
    For any $u\in \R$ with probability one the function $\lambda\mapsto \alpha_\lambda(u)$ is analytic and strictly increasing in $\lambda$ with $\alpha_0(u)=0$. 

    Moreover, $\alpha_\lambda(u)$ satisfies the following scale-invariance property: for any fixed $s\in \R$ 
\begin{align}\label{eq:scale_inv}
    \left(\alpha_\lambda(t), t\in (-\infty,\infty)\right) \ed  \left(\alpha_{e^{\beta s/4}\lambda}(t-s), t\in (-\infty,\infty)\right). 
\end{align}  
\end{enumerate}
\end{proposition}



% \benedek{Moved the statement about the flows to Section \ref{sec:additional}}




% As in proposition (\ref{prop:alphaSDE}) we have  $\alpha_\lambda(t)$ the unique strong solution of (\ref{eq:alphaSDE}):
% \begin{align}
% \label{eq:alphaSDEagain}
% 	d\alpha_\lambda =\lambda\tfrac{\beta}{4}e^{\frac{\beta}{4}u} du -\delta \sin(\alpha_\lambda)du+\Re[(e^{-i\alpha_\lambda}-1)dZ]
% \end{align}
% with $\lim_{t\to -\infty} \alpha_\lambda(t)=0$.
% Set $Q_\lambda(u)=e^{i  \alpha_\lambda(u)}$. By It\^o's formula we have, for any positive integer $k$:
% \begin{align}
% \label{eq:Qk_sde}
% 	&dQ^k= \left(\left(\cG_\delta(Q)\right)_k+i k \lambda f(u) Q^k \right)du-\frac{1}{2} i  k (Q-1) Q^{k-1} dZ+\frac{1}{2} i  k (Q-1) Q^k d\bar Z\\
%     \label{eq:SDE_1}
%         &d e^{-i k \lambda F+k^2 t}Q^k=\frac{k(k+\delta)}{2} e^{k^2 t-i k \lambda F(t)}Q^{k-1} dt+\frac{k(k-\delta)}{2}e^{-i k \lambda F(t)+k^2 t}Q^{k+1}dt \nonumber \\
%         &\qquad +\frac{k\left(-1+Q\right)}{2}   e^{ k^2 t -i \lambda F(t)} Q^{k-1}d \bar{Z}+ \frac{k(1+Q)}{2}   e^{+k^2 t-i k\lambda F(t) }Q^{k+1}  dZ  \\
%         \label{eq:SDE_2}
%         &d e^{-i k \lambda F}Q^k=\frac{k(k+\delta)}{2} e^{-i k \lambda F(t)}Q^{k-1}dt+\frac{k(k-\delta)}{2}e^{-i k \lambda F(t)}Q^{k+1}dt-k^2 e^{-i k \lambda F(t)}Q^k dt \nonumber \\
%       &\qquad+\frac{k\left(-1+Q\right)}{2}   e^{  -i \lambda F(t)} Q^{k-1}d \bar{Z}+ \frac{k(1+Q)}{2}   e^{-i k\lambda F(t) } Q^{k+1} dZ 
% \end{align}






\subsection{Palm measures}

The Palm measure of a stationary point process $\Xi$ on $\R$ is the process viewed from a `typical' point, with that point shifted to 0. The reduced Palm measure is the resulting process with 0 removed. 
For a precise definition, see Section 13 of \cite{DVJ2}  or \cite{BVBV_Palm}. 
Section 13 of \cite{DVJ2}  shows that for a stationary process $\Xi$ with a finite pair correlation function $\rho_{\Xi}^{(2)}(0,x)$, the density  of the (reduced) Palm measure exists and it is given by 
\[\rho_{Palm}^{(1)}(\lambda)=\frac{\rho_{\Xi}^{(2)}(0,x)}{\rho_{\Xi}^{(1)}(x)}.
\]



\begin{proposition}\label{prop:Palm}
Let $\beta>0$.   The reduced Palm measure of $\Sineb$ is given by $\HPbb$.
\end{proposition}
\begin{proof}
The proof of this statement is basically contained in \cite{BVBV_Palm}, although it is not stated there explicitly. 


Note that we can readily condition a size $n+1$ circular beta-ensemble to have a point at angle 0: we just have to set one of the angles to be 0 in the joint density function, and renormalize it. The resulting joint density is exactly $\textup{CJ}_{n,\beta,\beta/2}$ (see \eqref{eq:PDF_cjacobi}).
This shows that the Palm measure of $\textup{C}_{n+1,\beta}$ is $\textup{CJ}_{n,\beta,\beta/2}$, see Proposition 6 of \cite{BVBV_Palm} for a detailed proof and additional discussions. 
Lemma 4 of \cite{BVBV_Palm}  shows that the Palm measure of  $\Sineb$ can be obtained as the point process limit of $\textup{CJ}_{n,\beta,\beta/2}$ scaled with $n$ with an extra point added at 0. From \cite{LV_HP}  we know that the limit of this process (without the point at zero) is  $\HPbb$, which identifies $\HPbb$ as the reduced Palm measure of $\Sineb$.
\end{proof}



Proposition \ref{prop:Palm} together with the discussion before it leads to the following statement.
\begin{proposition}\label{prop:sine_HP}
    Fix $\beta>0$ and let $\rho_{\beta,\beta/2}^{(1)}$ be the density (one-point function) of the  $\HPbb$ process. Let $\rho_{\beta}^{(2)}$ denote the pair correlation (two-point function) of the  $\Sineb$ process. Then
    \begin{align}
        \rho_{\beta}^{(2)}(0,\lambda)=\frac{1}{2\pi} \rho_{\beta,\beta/2}^{(1)}(\lambda).\label{eq:sine_HP}
    \end{align}
\end{proposition}




\section{The density of $\HPb$ for $\delta>0$}
\label{sec:1point}

The main result of this section is the following theorem. 

\begin{theorem}[One-point function of $\HPb$]\label{thm:1p_HP}
 Let $\delta>0$ and $\beta>0$. 
Consider the strong solution $\alpha_\lambda(u)$ of \eqref{eq:alphaSDE} with initial condition \eqref{eq:SDE_initial}. Then the one-point function of the $\HPb$ process is given by 
    \begin{align}\label{eq:rho1_sum}
        \rho^{(1)}_{\beta,\delta}(\lambda)=\frac{1}{2\pi}+\frac{1}{\pi}\sum_{k=1}^\infty  \frac{(-\delta)^{\uparrow k}}{(1+\delta)^{\uparrow k}} E[\cos(k \alpha_\lambda(0)]. 
    \end{align}
\end{theorem}
By the bound \eqref{eq:coeff_bnd} the sum on the right side of \eqref{eq:rho1_sum} is absolutely convergent.


\begin{proof}[Proof of Theorem \ref{thm:sineb_corr}]
Using Theorem \ref{thm:1p_HP} with $\delta=\beta/2$ together with the identity \eqref{eq:sine_HP} of Proposition \ref{prop:sine_HP} yields the statement of Theorem \ref{thm:sineb_corr}.
\end{proof}

In order to  prove Theorem \ref{thm:1p_HP} we  first average out  $\Theta$ in the characterization \eqref{eq:HP} of the $\HPb$ process. 


% We start with a result on the one-point function of a point process generated as the result of a shooting problem. 
\begin{lemma}\label{lem:exp_count}
Suppose that $g : \R\mapsto \R$ is a continuous, strictly increasing deterministic function with $g(0)=0$. Fix $\delta>0$, and let  $\Theta\in[0,2\pi)$ be a random variable with density  \eqref{eq:Theta}. Consider the process
\[
\Xi=\{\lambda\in \R: g(\lambda)=\Theta \mod 2\pi\}.
\]
Then for $\lambda>0$ the expected number of points of $\Xi$ in $[0,\lambda]$ is given by 
\begin{align}\label{eq:Xi_counting}
E[\Xi[0,\lambda]]&=\frac{g(\lambda)}{2\pi}+\frac{1}{\pi}\sum_{k=1}^\infty  \frac{(-\delta)^{\uparrow k}}{(1+ \delta)^{\uparrow k}}\frac{\sin( k g(\lambda))}{k}.
\end{align}
% Note that if $\delta=n$ then the sum can be taken from $k=1$ to $n$.
\end{lemma}

\begin{proof}
Since $\lambda>0$, $g(0)=0$, and $g$ is continuous and strictly increasing, we have
\[
\Xi[0,\lambda]=\big|\{k\in \Z_{\ge 0}: \Theta+2k\pi\le g(\lambda) \}\big|=\lfloor\tfrac{g(\lambda)}{2\pi} \rfloor+\ind \left(g(\lambda)-2\pi \lfloor\tfrac{g(\lambda)}{2\pi}\rfloor\ge \Theta \right),
\]
and
\[
E[\Xi[0,\lambda]]=\lfloor\tfrac{g(\lambda)}{2\pi} \rfloor+P\left(g(\lambda)-2\pi \lfloor\tfrac{g(\lambda)}{2\pi}\rfloor\ge \Theta \right).
\]
Denote the probability density function of $\Theta$ by $h_\delta(\theta)$, the lemma will follow if we show that  for $0\le \theta<2\pi$ we have
\begin{align} \label{eq:CDF_h}
    \int_0^x  h_\delta(\theta) d\theta=\frac{x}{2\pi}+\frac1{\pi} \sum_{k=1}^\infty  \frac{(-\delta)^{\uparrow k}}{(1+\delta)^{\uparrow k}}\cdot \frac{\sin(k x)}{k}. 
\end{align}
The function $h_\delta:[0,2\pi)\to \R_+$ is bounded and real, hence it has a Fourier expansion 
\[
h_\delta(\theta)=\sum_{k\in \Z} a_k e^{i k \theta}
\]
with $a_{-k}=\bar a_k$. We will show that 
\[
a_k=\frac{1}{2\pi} \frac{(-\delta)^{\uparrow k}}{(1+\delta)^{\uparrow k}}, \qquad \text{for $k\ge 0$,}
\]
from this \eqref{eq:CDF_h} follows directly by integration. 
% $a_0=\frac{1}{2\pi}$ follows from the fact that $h_\delta$ is a probability density function. 
We have 
\begin{align*}
    (1-z)^a=\sum_{k=0}^\infty \frac{(-a)^{\uparrow k}}{k!} z^j,
\end{align*}
where this series is absolutely convergent for $|z|=1$ if $ a>0$. Since $\delta>0$, we have
\begin{align*}
     (1-e^{i \theta})^{\delta}(1-e^{-i \theta})^{ \delta}&= \sum_{m, n\ge 0}
\frac{(- \delta)^m (-\delta)^n}{m! n!} e^{i \theta (m-n)},
\end{align*}
where the sum is also absolutely convergent. For $k\ge 0$ the coefficient of $e^{i k \theta}$ is given by
\begin{align*}
    \sum_{n=0}^\infty \frac{(-\delta)^{\uparrow n+k} (-\delta)^{\uparrow n}}{(n+k)! n!}&=
    (-\bar \delta)^{\uparrow k}\sum_{n=0}^\infty  \frac{(- \delta+k)^{\uparrow n} (-\delta)^{\uparrow n}}{(n+k)! n!}\\
    &=(-\delta)^{\uparrow k} \frac{1}{k!} { }_2F_1(-\delta+k, -\delta, k+1,1)\\
    &=(- \delta)^{\uparrow k}  \frac{ \Gamma(1+2\delta)}{\Gamma(1+ \delta)\Gamma(1+\delta+k)},
\end{align*}
where the last two equations follow from the definition of the $_2F_1$ hypergeometric function and the Gauss formula \cite{andrews1999special}. From this it follows that for $k\ge 0$ the Fourier coefficient $a_k$ of the function $h_\delta$ is 
\begin{align*}
    a_k=\frac{1}{2\pi} \frac{\Gamma(1+\delta)^2}{\Gamma(1+2\delta)} \cdot  (-\delta)^{\uparrow k}  \frac{ \Gamma(1+2\delta)}{\Gamma(1+ \delta)\Gamma(1+\delta+k)}=\frac{1}{2\pi}\frac{ (- \delta)^{\uparrow k} }{(1+\delta)^{\uparrow k}},
\end{align*}
which is what we wanted to prove.
\end{proof}


\begin{definition}\label{def:Gd}
We define the operator $\cG_\delta:\R^{\ZZ_{+}}\to \R^{\ZZ_{+}}$ acting on sequences $(a_1,a_2,\dots)$ as 
\begin{align}\label{def:G_delta}
    \left(\cG_\delta(a)\right)_k=\frac{k(k+\delta)}{2}a_{k-1}-k^2 a_k+\frac{k(k-\delta)}{2} a_{k+1}, \qquad k\ge 1,
\end{align}
where we set $a_0$ to be 1. 
\end{definition}


The next claim follows directly from the definition of the Pochhammer symbol and partial summation. 
\begin{claim} \label{cl:partial_sum} For any integer $m\ge 1$ we have
    \begin{align}
          \sum_{k=1}^m \frac{(-\delta)^{\uparrow k}}{k (1+\delta)^{\uparrow k}}\left(\cG_\delta (a)\right)_k=\frac{\delta}{2}(a_1-1)-\frac{1}{2}(m-\delta) \frac{(-\delta)^{\uparrow m}}{ (1+\delta)^{\uparrow m}}(a_m-a_{m+1}).
    \end{align}
\end{claim}
% \begin{proof}
% This follows directly from the fact that for $k\ge 1$ we have
% \[
% \frac{k+1+\delta}{2}\cdot \frac{(-\delta)^{\uparrow k+1}}{(1+\delta)^{\uparrow k+1}}-k \cdot \frac{(-\delta)^{\uparrow k}}{ (1+\delta)^{\uparrow k}}+\frac{k-1-\delta}{2} \frac{(-\delta)^{\uparrow k-1}}{k (1+\delta)^{\uparrow k-1}}=0.
% \]
% % definition of $\coeff_{k,\delta}$ and $\cG_\delta (\cdot)_k$ (see  \eqref{eq:coeff}, \eqref{def:G_delta}) together with $$\frac{(\delta+1)^{\uparrow k+1}}{k+1+\delta} = (\delta+1)^{\uparrow k}.$$  
% \end{proof}



Consider the solution $\alpha_\lambda$ of the SDE \eqref{eq:alphaSDE}. It\^{o}'s formula shows that  the evolution of $e^{i k \alpha_\lambda(u)}$ for $\delta>0$, $k\ge 1$ is given by :  
\begin{align}\notag
 d e^{i k\alpha_\lambda(u)}&=\frac{k(k+\delta)}{2} e^{i (k-1)\alpha_\lambda(u)} du+(i k \lambda \tfrac{\beta}{4}e^{\frac{\beta}{4}u}-k^2) e^{i k \alpha_\lambda(u)} du+\frac{k(k-\delta)}{2} e^{i (k+1)\alpha_\lambda(u)}du  \\
 &\qquad -\frac{1}{2} i  k (e^{i \alpha_\lambda(u)}-1) e^{i(k-1)\alpha_\lambda(u)} dZ+\frac{1}{2} i  k (e^{i \alpha_\lambda(u)}-1) e^{i k \alpha_\lambda(u)} d\bar Z. \label{eq:QkSDE}
\end{align}
% Note that the same SDEs hold for $\alpha_{\lambda,\nu}(u)$ and $e^{i k\alpha_{\lambda,\nu}}(u)$ as well. \benedek{add connection to $\cG$} 
Note that the drift term in this SDE can be written as 
\[
 \left(\cG_\delta(e^{i j \alpha_\lambda(u)})\right)_k+i k \lambda \tfrac{\beta}{4}e^{\frac{\beta}{4}u} e^{i k \alpha_\lambda(u)},
\]
where the first term is the $k$th coordinate of the image of the operator $\cG_\delta$ acting on the sequence $e^{i j \alpha_\lambda(u)}, j\ge 1$.

We also record the fact that when  $\delta>0$ the SDE \eqref{eq:alphaSDE} becomes
\begin{align}
    \label{eq:alphaSDE11}
	d\alpha_\lambda = \lambda\tfrac{\beta}{4}e^{\frac{\beta}{4}u} du -\delta \sin(\alpha_\lambda(u) du+\Re[(e^{-i\alpha_\lambda(u)}-1)dZ].
\end{align}




\begin{proof}[Proof of Theorem \ref{thm:1p_HP}]
It is sufficient to prove the statement for $\lambda>0$, the other case follows by symmetry.  Recall the construction of $\HPb$ from \eqref{eq:HP} and the properties of $\alpha_\lambda$ from Proposition \ref{prop:alphaSDE}.
Since the process $\alpha_\lambda$ is independent of $\Theta$, and $\lambda\mapsto \alpha_\lambda(0)$ satisfies the  conditions required for $g(\lambda)$ in Lemma \ref{lem:exp_count}, 
we obtain
\begin{align}
  \pi  E[\HPb[0,\lambda]]&=\frac12 E[\alpha_\lambda(0)]+\sum_{k=1}^\infty \frac{(-\delta)^{\uparrow k}}{k (1+\delta)^{\uparrow k}}  \Im E[e^{i k \alpha_\lambda(0)}].
\end{align}
The sum on the right is absolutely convergent by the bound \eqref{eq:coeff_bnd}.

% Note that the coefficient of the noise coefficient is  bounded by a $k$-dependent constant. 

% If we could take expectations in \eqref{eq:alphaSDEagain} and \eqref{eq:Qk_sde}  then we would get
% \begin{align}\label{eq:ode_alpha}
%    \frac{d}{du} E[\alpha_\lambda(u)]&=\lambda f(u)-\delta \Im[E[Q_\lambda(u)]]\\
%    \frac{d}{du} E[Q^k_{\lambda}(u)]&=\left(\cG_\delta E[Q^{j}_\lambda(u)]\right)_k+ ik \lambda f(u) E[Q^k_\lambda(u)]\label{eq:ode_Q}
% \end{align}
% Integrating from $-\infty$ (noting the initial condition $\alpha_\lambda(-\infty)=0$), \emph{formally} \yahui{Can we say formally???}
% \begin{align}
%     \frac{1}{2}E[\alpha_\lambda(0)]+\sum_{k=1}^\infty \coeff_{k,\delta} \Im  E[e^{i k \alpha_\lambda(0)}]=&\frac{\lambda}{2}-\frac{\delta}{2} \Im \int_{-\infty}^0 E[Q_\lambda(u)]du+ \int_{-\infty}^0 \Im \sum_{k=1}^\infty \coeff_{k,\delta} \left(\cG_\delta E[Q_\lambda^j(u)]\right)_k du\notag\\
%     &
%     +\lambda \int_{-\infty}^0 \sum_{k=1}^\infty f(u) \coeff_{k,\delta} k  \Re[E[Q^k_\lambda(u)]] du
% \end{align}
% By Claim \ref{cl:partial_sum} (since the boundary term vanishes) we get
% \[
% \Im  \sum_{k=1}^\infty \coeff_{k,\delta} \left(\cG_\delta E[Q_\lambda^j(u)]\right)_k=\frac{\delta}{2}\Im (E[Q_\lambda(u)]-1)=\frac{\delta}{2} \Im E[Q_\lambda(u)].
% \]
% Thus we get
% \[
%  \pi E[\HPb[0,\lambda]]=\frac{\lambda}{2}+\lambda \int_{-\infty}^0 \sum_{k=1}^\infty f(u) k \coeff_{k,\delta}\Re[E[Q^k_\lambda(u)]] du.
% \]
% By the scale invariance of the $\alpha$ process we have
% \[
% E[Q_\lambda^k(u))]=E[Q^k_{\lambda e^{\beta u/4}}(0)]
% \]
% and introducing the change of variables $x=\lambda e^{\beta u/4}$ we get
% \begin{align}\label{eq:integrated_rho}
%  E[\HPb[0,\lambda]]=\frac{1}{2\pi}\int_0^\lambda \left(1+2\sum_{k=1}^\infty k \coeff_{k,\delta} E[\cos(k \alpha_x(0)]\right) dx,
% \end{align}
% identifying $\rho^{(1)}_{\beta,\delta}(\lambda)$ for $\lambda>0$ as the integrand on the right. 

 % It is not clear that we are allowed to integrate \eqref{eq:ode_alpha} and \eqref{eq:ode_Q} from $-\infty$ to 0. But the proof is valid on any compact interval, which together with the scaling properties of $\alpha_\lambda(u)$ (\ref{prop:alphaSDE}) allows  us to complete the proof.

By Proposition \ref{prop:alphaSDE} the function $\lambda\mapsto \alpha_\lambda(0)$ is monotone in $\lambda$, hence by the monotone convergence theorem 
\begin{align*}
  \pi   E[\HPb[0,\lambda]]&=\lim_{\eps\to 0^+} \pi E[\HPb[0,\lambda]-\HPb[0,\eps  ]]\\
  &=\lim_{\eps\to 0^+} \frac12 E[\alpha_\lambda(0)-\alpha_{\eps}(0)]+\sum_{k=1}^\infty \frac{(-\delta)^{\uparrow k}}{k (1+\delta)^{\uparrow k}} \Im E[e^{i k \alpha_\lambda(0)}-e^{i k \alpha_{\eps}(0)}].
\end{align*}
Choose $\eps=\lambda e^{\beta \nu/4}$ with $\nu<0$. Then by \eqref{eq:scale_inv} of Proposition \ref{prop:alphaSDE} we have $\alpha_{\eps}(0)\ed \alpha_\lambda(\nu)$, and hence
\begin{align*}
& \frac12 E[\alpha_\lambda(0)-\alpha_{\eps}(0)]+\sum_{k=1}^\infty \frac{(-\delta)^{\uparrow k}}{k (1+\delta)^{\uparrow k}}\Im E[e^{i k \alpha_\lambda(0)}-e^{i k \alpha_{\eps}(0)}]\\
&\qquad\qquad \qquad =\frac12 E[\alpha_\lambda(0)-\alpha_{\lambda}(\nu)]+\sum_{k=1}^\infty \frac{(-\delta)^{\uparrow k}}{k (1+\delta)^{\uparrow k}}\Im E[e^{i k \alpha_\lambda(0)}-e^{i k \alpha_{\lambda}(\nu)}]\\
&\qquad \qquad\qquad =\frac12 E[\alpha_\lambda(0)-\alpha_{\lambda}(\nu)]+\lim_{m\to \infty}\sum_{k=1}^m \frac{(-\delta)^{\uparrow k}}{k (1+\delta)^{\uparrow k}}\Im E[e^{i k \alpha_\lambda(0)}-e^{i k \alpha_{\lambda}(\nu)}].
\end{align*}
For a fixed $\lambda>0$, $k\ge 1$, $\nu<0$ the stochastic differential equations  \eqref{eq:QkSDE} and \eqref{eq:alphaSDE11}  have uniformly bounded coefficients for $u\in [\nu,0]$, so we can express $E[e^{i k \alpha_\lambda(0)}-e^{i k \alpha_{\lambda}(\nu)}]$ and $E[\alpha_\lambda(0)-\alpha_\lambda(\nu)]$ as the integrated expected drifts on $[\nu,0]$ from the respective  equations.

Fixing $m\ge 1$, and using the notation \eqref{def:G_delta}, we get
\begin{align*}
    &\frac12 E[\alpha_\lambda(0)-\alpha_{\lambda}(\nu)]+\sum_{k=1}^m \frac{(-\delta)^{\uparrow k}}{k (1+\delta)^{\uparrow k}}\Im E[e^{i k \alpha_\lambda(0)}-e^{i k \alpha_{\lambda}(\nu)}]\\
    &\qquad =\frac{\lambda}{2}(1-e^{\beta \nu/4})-\frac{\delta}{2} \Im \int_{\nu}^0 E[e^{i \alpha_\lambda(u)}]du+\Im  \int_{\nu}^0 \sum_{k=1}^m \frac{(-\delta)^{\uparrow k}}{k (1+\delta)^{\uparrow k}} \left(\cG_\delta E[e^{i j \alpha_\lambda(u)}]\right)_k du\\
    &\qquad\qquad
    +\lambda \int_{\nu}^0 \sum_{k=1}^m \frac{\beta}{4}e^{\frac{\beta}{4}u} \frac{(-\delta)^{\uparrow k}}{ (1+\delta)^{\uparrow k}}  \Re[E[e^{i k \alpha_\lambda(u)}]] du.
\end{align*}
Using Claim \ref{cl:partial_sum} for the sequence $a_j=E[e^{i j \alpha_\lambda(u)}]$ we get for any $m\ge 1$,
\begin{align}\notag  &-\frac{\delta}{2} \Im E[e^{i \alpha_\lambda(u)}]+\Im   \sum_{k=1}^m  \frac{(-\delta)^{\uparrow k}}{k (1+\delta)^{\uparrow k}}\left(\cG_\delta E[e^{i j \alpha_\lambda(u)}]\right)_k\\&\notag \qquad  = -\frac{\delta}{2} \Im E[e^{i \alpha_\lambda(u)}]+\frac{\delta}{2}\Im[E[e^{i \alpha_\lambda(u)}-1]]-\frac12 (m-\delta)\frac{(-\delta)^{\uparrow m}}{ (1+\delta)^{\uparrow m}}\Im E[e^{i m \alpha_\lambda(u)}-e^{i (m+1) \alpha_\lambda(u)}]\\
   &\qquad  =-\frac12 (m-\delta)\frac{(-\delta)^{\uparrow m}}{ (1+\delta)^{\uparrow m}}\Im E[e^{i m \alpha_\lambda(u)}-e^{i (m+1) \alpha_\lambda(u)}].\label{eq:error}
\end{align}
Note that by the bound  \eqref{eq:coeff_bnd}, the term \eqref{eq:error} is uniformly bounded by $c m^{-2\delta}$, hence for a fixed $\nu<0$ this term will vanish when integrated on $[\nu,0]$ and then $m$ taken to $\infty$.  This leads to 
\begin{align*}
&   \frac12 E[\alpha_\lambda(0)-\alpha_{\lambda}(\nu)]+\lim_{m\to \infty}\sum_{k=1}^m \frac{(-\delta)^{\uparrow k}}{k (1+\delta)^{\uparrow k}}\Im E[e^{i k \alpha_\lambda(0)}-e^{i k \alpha_{\lambda}(\nu)}]\\
&\qquad\qquad =\frac{\lambda}{2}(1-e^{\beta \nu/4}) +\lambda \lim_{m\to \infty}\int_{\nu}^0 \sum_{k=1}^m \frac{\beta}{4}e^{\frac{\beta}{4}u} \frac{(-\delta)^{\uparrow k}}{ (1+\delta)^{\uparrow k}}  \Re[E[e^{i k \alpha_\lambda(u)}]] du\\
&\qquad\qquad =\frac{\lambda}{2}(1-e^{\beta \nu/4}) + \int_\nu^0 \sum_{k=1}^\infty  \frac{\beta}{4}e^{\frac{\beta}{4}u} \frac{(-\delta)^{\uparrow k}}{ (1+\delta)^{\uparrow k}}  E[\cos( k \alpha_\lambda(u))]] du\\
&\qquad\qquad =\frac{\lambda}{2}(1-e^{\beta \nu/4}) + \int_\eps^\lambda \sum_{k=1}^\infty  \frac{\beta}{4}e^{\frac{\beta}{4}u} \frac{(-\delta)^{\uparrow k}}{ (1+\delta)^{\uparrow k}}  E[\cos( k \alpha_x(0))]] dx.
\end{align*}
The  bound \eqref{eq:coeff_bnd} shows that the series is absolutely convergent, while $\eps=\lambda\frac{\beta}{4}e^{\frac{\beta}{4} \nu}$ together with the scale invariance \eqref{eq:scale_inv} justifies the last step.  Letting $\eps\to 0$,  the bounded convergence theorem yields \[
 \pi E[\HPb[0,\lambda]]=\frac{\lambda}{2}+\int_{0}^\lambda \sum_{k=1}^\infty    \frac{(-\delta)^{\uparrow k}}{ (1+\delta)^{\uparrow k}}  E[\cos(k \alpha_\lambda(u)] du,
\] and the statement of Theorem \ref{thm:1p_HP} follows.
\end{proof}



\section{Large $\lambda$ bounds}\label{sec:large_lambda}

The goal of this section is to prove Theorem \ref{thm:large_lambda}. The main ingredient is the following lemma. 
\begin{lemma}\label{lem:E_alpha}
Fix $\beta>0$ and let $\alpha_\lambda(t)$ be the strong solution of \eqref{eq:alphaSDE_b}. There is a positive finite constant $c$ only depending on $\beta$ so that the following bounds hold for $\lambda\ge 2$. 
\begin{align}\label{eq:eialpha_1}
   \big|E[e^{i \alpha_\lambda(0)}]\big|
   &\le c \left(\lambda^{-1} + \lambda^{-4/\beta}\right), \\[0.5em]
   \big|E[e^{i k \alpha_\lambda(0)}]\big|
   &\leq c  \Big(k^2 \lambda^{-1}+   \log(\lambda)\,\lambda^{-4k^2/\beta} \cdot \ind(4k^2\le \beta)\Big), \qquad k\ge 2.\label{eq:eialpha_k}
\end{align}
\end{lemma}

We first show how Theorem \ref{thm:large_lambda} follows from the lemma.
\begin{proof}[Proof of Theorem \ref{thm:large_lambda}]
In the arguments below $c$ is always a positive constant depending only on $\beta$, but its value might change from line to line. 

From Theorem \ref{thm:sineb_corr} and the bound \eqref{eq:coeff_bnd} we obtain
\begin{align}\label{eq:trunc_bnd}
    \left| \rho_\beta^{(2)}(0,\lambda)-\frac{1}{4\pi^2}\right|\le c\sum_{k=1}^\infty  k^{-1-\beta} \left|E[e^{i k \alpha_\lambda(0)}]\right|.
\end{align}
When $\beta > 2$ we use \eqref{eq:eialpha_1} and \eqref{eq:eialpha_k} in each term of \eqref{eq:trunc_bnd} to get
\begin{align*}
    \left| \rho_\beta^{(2)}(0,\lambda)-\frac{1}{4\pi^2}\right|
    &\leq c \big(\lambda^{-1}+\lambda^{-4/\beta}\big)
      + c \sum_{k=2}^\infty  k^{1-\beta} \lambda^{-1}+c \!\!\sum_{2\le k\le \sqrt{\beta/4}} k^{-1-\beta}\log(\lambda) \lambda^{-4k^2/\beta}\\
&\le c(\lambda^{-1}+\lambda^{-4/\beta}).
\end{align*}
When $0 < \beta < 2$ we bound $\left|E[e^{i k \alpha_\lambda(0)}]\right|$ by 1 for $k>\lambda^{1/2}$, and use \eqref{eq:eialpha_1}-\eqref{eq:eialpha_k} otherwise. Note that in this case $4k^2>\beta$ holds for all $k$.
\begin{align*}
    \left| \rho_\beta^{(2)}(0,\lambda)-\frac{1}{4\pi^2}\right|
     &\le c\big(\lambda^{-1}+\lambda^{-4/\beta}\big)
      + c\sum_{2\le k\le \lambda^{1/2}}
  k^{1-\beta} \lambda^{-1}+ c \sum_{k> \lambda^{1/2}} k^{-1-\beta}\\
    &\le c \lambda^{-\beta/2}.
\end{align*}
Finally, when $\beta=2$, the sum on the right of \eqref{eq:twopoint_cor} only contains a single term, $E[\cos(\alpha_\lambda(0))]$, hence from \eqref{eq:eialpha_1} we obtain the desired upper bound. Of course, for $\beta=2$ we also have the exact form of $\rho_2^{(2)}$ from \eqref{eq:sine_kernel}, which would yield a more precise bound.
\end{proof}


Now we turn to the proof of Lemma \ref{lem:E_alpha}. 
\begin{proof}[Proof of Lemma \ref{lem:E_alpha}]
Fix $k\in \Z_+$. From \eqref{eq:alphaSDE_b} and It\^o's formula we see that the evolution of the process $e^{i k \alpha_\lambda(u) - i k \lambda e^{\frac{\beta}{4}u} + k^2 u}$ is governed by the following SDE:
\begin{align}
    \label{eq:SDE_1}
    d e^{i k \alpha_\lambda(u) - i k \lambda e^{\frac{\beta}{4}u} + k^2 u}
    &= \frac{k(k+\frac{\beta}{2})}{2}\, e^{k^2 u - i k \lambda e^{\frac{\beta}{4}u}} e^{i (k-1)\alpha_\lambda(u)}\, du
     + \frac{k(k-\frac{\beta}{2})}{2}\, e^{k^2 u - i k \lambda e^{\frac{\beta}{4}u}} e^{i (k+1)\alpha_\lambda(u)}\, du \nonumber \\
    &\quad + \frac{k(-1+e^{i\alpha_\lambda(u)})}{2}\,
        e^{i(k-1)\alpha_\lambda(u) + k^2 u - i\lambda e^{\frac{\beta}{4}u}}\, d\bar{Z} \nonumber \\
    &\quad + \frac{k(1+e^{i\alpha_\lambda(u)})}{2}\,
        e^{i(k+1)\alpha_\lambda(u) + k^2 u - i k\lambda e^{\frac{\beta}{4}u}}\, dZ.
\end{align}
Note that the coefficient of each term in this SDE is bounded by a $k$-dependent constant times $e^{k^2 u}$. Hence by $\lim\limits_{u\to -\infty} \alpha_\lambda(u)=0$ and the bounded convergence theorem we can compute $E[e^{i k \alpha_\lambda(0)}]=E[e^{i k \alpha_\lambda(0) - i k \lambda e^{\frac{\beta}{4}\cdot 0} + k^2\cdot 0}]$ by integrating the expected values of the drift terms in \eqref{eq:SDE_1} on $(-\infty,0]$. 
\begin{align}\notag
   E[e^{i k \alpha_\lambda(0)}]&=\int_{-\infty}^0  \frac{k(k+\frac{\beta}{2})}{2}\, e^{k^2 u - i k \lambda e^{\frac{\beta}{4}u}} E[e^{i (k-1)\alpha_\lambda(u)}]\, du\\
   &\quad+\int_{-\infty}^0\frac{k(k-\frac{\beta}{2})}{2}\, e^{k^2 u - i k \lambda e^{\frac{\beta}{4}u}} E[e^{i (k+1)\alpha_\lambda(u)}]\, du
   \label{eq:E_qk}
\end{align}
We estimate $ | E[e^{i k \alpha_\lambda(0)}]|$ by estimating the absolute values of the integrals in \eqref{eq:E_qk}. 
Let $T_\lambda:=-\frac{4 \log \lambda}{\beta}< 0$. On the interval $(-\infty,T_\lambda]$, we bound the integrals by taking the absolute values inside together with $|E[e^{i j \alpha_\lambda(u)}]|\le 1$ to obtain
\begin{align}\notag
& \left| \int_{-\infty}^{T_\lambda}  \frac{k(k+\frac{\beta}{2})}{2}\, e^{k^2 u - i k \lambda e^{\frac{\beta}{4}u}} E[e^{i (k-1)\alpha_\lambda(u)}]\, du  +\int_{-\infty}^{T_\lambda}\frac{k(k-\frac{\beta}{2})}{2}\, e^{k^2 u - i k \lambda e^{\frac{\beta}{4}u}} E[e^{i (k+1)\alpha_\lambda(u)}]\, du\right|\\
&\qquad \le \frac{k(k+\tfrac{\beta}{2})+k|k-\tfrac{\beta}{2}|}{2}\int_{-\infty}^{T_\lambda} e^{k^2 u} du=\frac{k(k+\tfrac{\beta}{2})+k|k-\tfrac{\beta}{2}|}{2k^2} \lambda^{-\frac{4k^2}{\beta}}\le (\tfrac{\beta}{2}+1) \lambda^{-\frac{4k^2}{\beta}}.\label{eq:term_97}
\end{align}
In order to control the integrals in \eqref{eq:E_qk} on  $[T_\lambda,0]$, we need more delicate estimates. 
For $T_\lambda\le u\le 0$ we set
\[
    \Lambda_k(u)
    := \int_{T_\lambda}^u e^{-i \lambda e^{\frac{\beta}{4}s}+ k^2 s}\, ds
    = \frac{4}{\beta}\, \lambda^{-4k^2/\beta}
        \int_1^{\lambda e^{\beta u/4}} e^{-i u} u^{\frac{4k^2}{\beta}-1}\, du.
\]
Integration by parts shows that  
\begin{align}\label{eq:llasym_ibp}
          \left|\int_1^v e^{-i u} u^{a-1}\,du\right|
    \;\leq\; 2 (v^{a-1}+\ind(a<1)), \qquad \text{for all $a>0$, $v>1$.}
\end{align}
This implies the bound
\begin{align}\label{eq:bd_Lambdak}
    |\Lambda_k(u)|
    \;\leq\; \frac{8}{\beta}\Big(\lambda^{-4k^2/\beta}\cdot \ind(4k^2<\beta)
    + \lambda^{-1}\, e^{\frac{\beta}{4}u(\frac{4k^2}{\beta}-1)}\Big), \qquad T_\lambda\le u\le 0.
\end{align}
Note that for $k=1$ we have
\begin{align}\label{eq:Lambda1_99}
\int_{T_\lambda}^0  \frac{k(k+\frac{\beta}{2})}{2}\, e^{k^2 u - i k \lambda e^{\frac{\beta}{4}u}} E[e^{i (k-1)\alpha_\lambda(u)}]\, du=\tfrac{1+\frac{\beta}{2}}{2} \Lambda_1(0), 
\end{align}
which can be bounded by using \eqref{eq:bd_Lambdak}.

% When $k=1$, the first integral in \eqref{eq:E_qk} becomes deterministic, and we can bound it on $[T_\lambda,0]$ as follows
% \[
% \left|\int_{T_\lambda}^0 \frac{(1+\frac{\beta}{2})}{2}\, e^{ u - i \lambda e^{\frac{\beta}{4}u}} \, du\right|\le \tfrac{2+\beta}{2} |\Lambda_1(0)|\le \tfrac{16+8\beta}{\beta}(\lambda^{-4/\beta}+\lambda^{-1}).
% \]
Fix  $k\ge 2$.  Applying Itô’s formula for the process $\Lambda_k(u) \cdot  e^{i (k-1)\alpha_\lambda(u)-i \lambda (k-1)e^{\frac{\beta}{4}u}}$ yields
% \[
% \Lambda_k(u) \cdot  e^{i (k-1)\alpha_\lambda(u)-i \lambda (k-1)e^{\frac{\beta}{4}u}}, \qquad u\in [T_\lambda, 0]
% \]
% to get
\begin{align}\label{eq:stoch_part_int}
  \Lambda_k(0)   e^{i (k-1)\alpha_\lambda(0)}=\int_{T_\lambda}^0  e^{k^2 u - i k \lambda e^{\frac{\beta}{4}u}} e^{i (k-1)\alpha_\lambda(u)}\, du+\int_{T_\lambda}^0 \Lambda_k(u) \,d\!\left[e^{i (k-1)\alpha_\lambda(u)-i \lambda (k-1)e^{\frac{\beta}{4}u}}\right]. 
\end{align}
% When $k=1$, the last term in \eqref{eq:stoch_part_int} is 0, and we get
% \begin{align}\label{eq:term_95}
% \left|\int_{T_\lambda}^0  e^{ u - i \lambda e^{\frac{\beta}{4}u}} \, du\right|=  |\Lambda_1(0)|\le \tfrac{8}{\beta}(\lambda^{-4/\beta}+\lambda^{-1}).
% \end{align}
It\^o's formula shows that in $d[e^{i (k-1) \alpha_\lambda - i \lambda (k-1) e^{\frac{\beta}{4}u}}]$ on the interval $[T_\lambda,0]$ the drift terms are uniformly bounded by $(2+\tfrac{\beta}{2})(k-1)^2$ and the coefficients of $dZ$ and $d\bar Z$ are both bounded by  $k-1$.
% \begin{align}
% \label{eq:SDE_2}
%     de^{i j \alpha_\lambda - i \lambda j e^{\frac{\beta}{4}u}}
%     &= \frac{j(j+\frac{\beta}{2})}{2}\, e^{i (j-1)\alpha_\lambda(u) - i j \lambda e^{\frac{\beta}{4}u}}\, du
%      + \frac{j(j-\frac{\beta}{2})}{2}\, e^{i (j+1)\alpha_\lambda(u) - i j \lambda e^{\frac{\beta}{4}u}}\, du \nonumber \\
%     &\quad - j^2 e^{i j \alpha_\lambda(u) - i j \lambda e^{\frac{\beta}{4}u}}\, du
%      + \frac{j(-1+e^{i \alpha_\lambda(u)})}{2}\,
%        e^{i (j-1)\alpha_\lambda(u) - i \lambda e^{\frac{\beta}{4}u}}\, d\bar{Z} \nonumber \\
%     &\quad + \frac{j(1+e^{i \alpha_\lambda(u)})}{2}\,
%        e^{i (j+1)\alpha_\lambda(u) - i j\lambda e^{\frac{\beta}{4}u}}\, dZ.
% \end{align}
% The drift terms in \eqref{eq:SDE_2} are bounded by $(2+\tfrac{\beta}{2})j^2$ on $[T_\lambda,0]$ and the coefficients of $dZ$ and $d\bar Z$ are both bounded by  $j$,
Hence after taking expectations in \eqref{eq:stoch_part_int} and using the bound \eqref{eq:bd_Lambdak} we obtain
\begin{align}\notag
   & \left|E \int_{T_\lambda}^0  e^{k^2 u - i k \lambda e^{\frac{\beta}{4}u}} e^{i (k-1)\alpha_\lambda(u)}\, du\right|\le |\Lambda_k(0)|+(2+\tfrac{\beta}{2})(k-1)^2\int_{T_\lambda}^0 |\Lambda_k(u)|du\\\notag
    &\qquad \le c (\lambda^{-4k^2/\beta} \ind(4k^2<\beta)+\lambda^{-1})+c (k-1)^2 \int_{T_\lambda}^0 (\lambda^{-4k^2/\beta} \ind(4k^2<\beta)
    + \lambda^{-1}\, e^{\frac{\beta}{4}u(\frac{4k^2}{\beta}-1)}\Big) du\\
    &\qquad\le  c \left(\lambda^{-4k^2/\beta} \log \lambda \cdot \ind(4k^2<\beta)+\lambda^{-1}\right)\label{eq:term_98}.
\end{align}
Here, and in the rest of this proof, $c$ is a constant only depending on $\beta$, and possibly changing from line to line. Using the same approach, we also get the following bound for any $k\ge 1$:
\begin{align}
       \left|E \int_{T_\lambda}^0  e^{k^2 u - i k \lambda e^{\frac{\beta}{4}u}} e^{i (k+1)\alpha_\lambda(u)}\, du\right|\le c \left(\lambda^{-4k^2/\beta} \log \lambda \cdot \ind(4k^2<\beta)+\lambda^{-1}\right)\label{eq:term_99}.
\end{align}
Combining our bounds, we obtain \eqref{eq:eialpha_k}. \\
To prove \eqref{eq:eialpha_1} we return to the intermediate estimate
\begin{align}
    |E[e^{i \alpha_\lambda(0)}]|&\le c (\lambda^{-\frac{4}{\beta}}+\lambda^{-1})+c \left|\int_{T_\lambda}^0 e^{u - i \lambda e^{\frac{\beta}{4}u}} E[e^{i 2\alpha_\lambda(u)}]\, du\right|
    % \\&\le c (\lambda^{-\frac{4}{\beta}}+\lambda^{-1})+c \int_{T_\lambda}^0 e^{u} \left|E[e^{i 2\alpha_\lambda(u)}]\right|\, du
\end{align}
which follows from \eqref{eq:E_qk}, \eqref{eq:term_97},  \eqref{eq:bd_Lambdak}, and \eqref{eq:Lambda1_99}.

The scale invariant property \eqref{eq:scale_inv}  of the $\alpha_\lambda$ diffusion together with the bound \eqref{eq:eialpha_k} for $k=2$ (with $\lambda^{-16/\beta} \log\lambda \cdot \ind(16\le \beta)$ replaced with $\lambda^{-8/\beta}$) leads to 
\begin{align*}
    \left|E[e^{i 2\alpha_\lambda(u)}]\right|= \left|E[\exp(i 2\alpha_{\lambda  e^{\frac{\beta}{4}u}}(0))]\right|\le c \left((\lambda e^{\frac{\beta}{4}u})^{-8/\beta}+(\lambda e^{\frac{\beta}{4}u})^{-1}\right).
\end{align*}
From this, we  obtain
\begin{align*}
  \left|\int_{T_\lambda}^0 e^{u - i \lambda e^{\frac{\beta}{4}u}} E[e^{i 2\alpha_\lambda(u)}]\, du\right|\le   \int_{T_\lambda}^0 e^{u} \left|E[e^{i 2\alpha_\lambda(u)}]\right|\, du\le c (\lambda^{-4/\beta}+\lambda^{-1}),
\end{align*}
which completes the proof of \eqref{eq:eialpha_1}.
\end{proof}





\section{Results for $\delta=n$}\label{sec:delta_n}

Our goal in this section is to prove Theorem \ref{thm:sineb_series} and  Corollary \ref{cor:small_lambda}, together with  similar results for $\rho_{\beta,n}^{(1)}(\lambda)$ with $n\in \Z_+$.\\
Fix $\beta>0$ and $\delta=n\in \Z_+$. In this case the infinite sum in Theorem \ref{thm:1p_HP} becomes a finite sum with $n$ terms, and the coefficients can be simplified as well. We have 
\begin{align}\label{eq:rho_HP_n}
      \rho^{(1)}_{\beta,n}(\lambda)=\frac{1}{2\pi}+\frac{1}{\pi}\sum_{k=1}^n (-1)^k  \frac{\binom{2n}{n+k}}{\binom{2n}{n}} \Re E[e^{i k \alpha_\lambda(0)}], 
\end{align}
with $\alpha_\lambda(u)$ the strong solution of \eqref{eq:alphaSDE} with initial condition \eqref{eq:SDE_initial} and $\delta=n$. 
Let
 \begin{align}
    \mathbf{q}(\lambda)=[q_1(\lambda), \dots, q_n(\lambda)]^T, \qquad q_k(\lambda)=E[e^{i k\alpha_\lambda(0)}].
\end{align}   
(We do not explicitly denote the dependence on $\beta$ and $n$.)
Recall the definition of $\mathbf{A}_n, \mathbf{B}_n, \mathbf{e}_n$  from \eqref{eq:AB}, \eqref{eq:vectors}.
Define $\mathbf{s}_n\in \CC^n , n\ge 0$
as
\begin{align}\label{eq:q_coeff_1}
    \mathbf{s}_0&=-\frac{n+1}{2}\mathbf{A}_n^{-1}\mathbf{e}_n, \qquad \mathbf{s}_k=i \left(k \mathbf{I}-\tfrac{4}{\beta} \mathbf{A}_n\right)^{-1} \mathbf{B}_n \mathbf{s}_{k-1}, \qquad k\ge 1.
\end{align}
Note that this agrees with the previous definition \eqref{eq:q_coeff} if $\beta=2n$, but we will need this sequence now for general $\beta>0$. From Lemma \ref{lem:An_eval} below, it will follow that $\mathbf{s}_k, k\ge 0$ are well defined. \\
We also introduce the notation
\begin{align}\label{eq:f}
    \mathbf{f}_n=[1,1,\dots,1]^T\in \R^n.
\end{align}
The main results of this section are the following theorem and its corollary. 

\begin{theorem}\label{thm:even_beta}
 The vector-valued function $\mathbf{q}(\lambda)$ satisfies the ODE
 \begin{align}\label{eq:ode_q_vec}
     \frac{\beta}{4} \lambda \mathbf{q}'(\lambda) = (i \frac{\beta}{4}\lambda  \mathbf{B}_n   +\mathbf{A}_n) \mathbf{q}(\lambda)+\frac{n+1}{2} \mathbf{e}_n, \qquad \mathbf{q}(0)=\mathbf{f}_n.
\end{align}
The functions $\lambda\mapsto q_k(\lambda)$ can be extended to $\CC$ as entire functions  so that \eqref{eq:ode_q_vec} holds for all $\lambda\in \CC$. With this extension, we have for all $\lambda\in \CC$
\begin{align}\label{eq:q_expansion}
    \mathbf{q}(\lambda)=\sum_{j=0}^\infty \mathbf{s}_j \lambda^j.
\end{align} 
\end{theorem}


\begin{corollary}
\label{cor:int_case_rho}
   With the notations introduced above, the density  of the $\operatorname{HP}_{\beta, n}$ process is given by
\begin{align}\label{eq:int_case_rho}
    \rho^{(1)}_{\beta,n}(\lambda) =\frac{1}{2\pi}\left(1+2\;\mathbf{v}_n^T\Re \mathbf{q}(\lambda)\right)=\frac{1}{\pi} \sum_{j=1}^\infty \mathbf{v}_n\cdot \mathbf{s}_{2j} \lambda^{2j}.
\end{align}
\end{corollary}
Using Proposition \ref{prop:sine_HP} with $\beta=2n$ gives the analogue statements for $\rho_{2n}^{(2)}(0,\lambda)$. Theorem \ref{thm:sineb_series}  follows from Corollary \ref{cor:int_case_rho}, and Theorem \ref{thm:even_beta} leads to a description of $\rho_{2n}^{(2)}(0,\lambda)$ in terms of the ODE \eqref{eq:ode_q_vec} (with $\beta=2n$). \\
We also identify the behavior of $\rho_{\beta,n}^{(1)}(\lambda)$ near 0. 
\begin{proposition}
\label{prop:small_lambda_asym}
    For $\lambda$ near 0, we have 
       \begin{align}\label{eq:small_lambda_HP}
        \rho_{\beta,n}^{(1)}(\lambda) =\frac{1}{2\pi} 
      \binom{2n}{n}^{-1}\frac{\left(\tfrac{\beta}{2}\right)^{2n}}{\left(1+\frac{\beta}{2}\right)^{\uparrow 2n}}   \lambda^{2n}  
        % \frac{1}{2\pi} \binom{2n}{n}^{-1} \prod_{j=1}^{2n}\left(1+\tfrac{4}{\beta}\Big(n-\tfrac{j-1}{2}\Big)\right)^{-1} \lambda^{2n}
        + \mathcal{O}(\lambda^{2n+2}).
    \end{align} 
\end{proposition}

When $\beta=2n$, we can readily check that 
\[
 \binom{2n}{n}^{-1}\frac{\left(\tfrac{\beta}{2}\right)^{2n}}{\left(1+\frac{\beta}{2}\right)^{\uparrow 2n}} =\frac{n^{2n} (n!)^3 }{(2n!)(3n!) },
\]
which immediately implies Corollary \ref{cor:small_lambda}. 












\begin{proof}[Proof of Theorem \ref{thm:even_beta}]
For a fixed $\nu<0$ consider the unique strong solution $\alpha_{\lambda,\nu}(u)$ of \eqref{eq:alphaSDE} on $[\nu,\infty)$ with initial condition $\alpha_{\lambda,\nu}(\nu)=0$ and $\delta=n$. Set 
\begin{align}
    q_{k,\nu}(u,\lambda)=E[e^{i k \alpha_{\lambda,\nu}(u)}].
\end{align}
Fix $\lambda\in \R$. By Proposition \ref{prop:alphaSDE} we have  $\lim\limits_{\nu\to -\infty} \alpha_{\lambda,\nu}(u)=\alpha_\lambda(u)$ almost surely, hence by the bounded convergence theorem we have 
\begin{align}
    \lim_{\nu\to -\infty}  q_{k,\nu}(0,\lambda)=q_k(\lambda). 
\end{align}
The process $e^{i k \alpha_{\lambda,\nu}(u)}$ satisfies the same SDE as $e^{i k \alpha_\lambda(u)}$, see \eqref{eq:QkSDE}. For a fixed $\lambda$, the coefficients of this SDE are bounded on $[\nu,0]$, hence the Brownian terms do not contribute when we take expectations in the integrated form of the equation. 
% It\^o's formula gives  
% \begin{align}
%         \label{eq:Q_sde}
%     dQ_\nu(u,\lambda)=&\frac{\delta+1}{2}du+(i \lambda f(u)-1) Q_\nu du+\frac{1}{2}(1-\delta) Q_\nu^2 du\\
%     &\notag \qquad-\frac{1}{2} (Q_\nu-1) (i db_1 Q_\nu-i db_1+db_2 Q_\nu+db_2)\\
%     dQ_\nu^k(u,\lambda)=&\frac{k}{2} (\delta +k)Q_\nu^{k-1}+(-k^2 + i k\lambda f_\beta(u))Q_\nu^{k}+\frac{1}{2} k (k-\delta )Q_\nu^{k+1}du\nonumber\\
%    &+\frac{1}{2} i k Q_\nu^{k-1} (1-Q_\nu)^2 db_1+\frac{1}{2} k Q_\nu^{k-1} \left(Q_\nu^2-1\right)db_2
% \end{align}
% The coefficients of these SDEs are bounded for $\nu\le u\le t$ (for any fixed $t>\nu$), hence we can take expectations to get 
This leads to the following ODE system for $u\in [\nu,\infty)$.
\begin{align}\label{eq:ODEq1}
    \partial_u q_{1,\nu}&=\frac{n+1}{2}+(i \lambda \tfrac{\beta}{4}e^{\frac{\beta}{4}u}-1) q_{1,\nu}+\frac{1}{2}(1-n)q_{2,\nu},\\\label{eq:ODEq2}
    \partial_u q_{k,\nu}&=\frac{k}{2} (n +k) q_{k-1,\nu}+(-k^2 + i k \lambda \tfrac{\beta}{4}e^{\frac{\beta}{4}u})q_{k,\nu}+\frac{1}{2} k (k-n )q_{k+1,\nu}, \qquad 2\le k\le n,\\[3pt]
   &q_{k,\nu}(\nu,\lambda)=1, \qquad 1\le k\le n. 
\end{align}
We did not denote the dependence on $u,\lambda$ in \eqref{eq:ODEq1} and \eqref{eq:ODEq2}. 
Since the coefficient 
of $q_{n+1,\nu}$ in \eqref{eq:ODEq2} is 0,   we obtained a closed ODE system.
Using the notation
\[
\mathbf{q}_{\nu}(u,\lambda)=[q_{1,\nu}(u,\lambda), \dots, q_{n,\nu}(u,\lambda)]^T,
\]
we can write it as 
\begin{align}\label{eq:q_nu_ODE}
     \partial_u \mathbf{q}_\nu(u,\lambda) &= (i \tfrac{\beta}{4}e^{\frac{\beta}{4}u} \mathbf{B}_n \lambda  +\mathbf{A}_n) \mathbf{q}_\nu(u,\lambda) +\frac{n+1}{2}\mathbf{e}_n, \qquad
     \mathbf{q}_{\nu}(\nu,\lambda) = \mathbf{f}_n.
\end{align}
For a fixed $\lambda\in \CC$ the coefficients of this ODE system are bounded for $u\in [\nu,0]$, hence there is unique solution on $[\nu,0]$ for each $\lambda\in \CC$, and this solution is analytic in $\lambda$ for a fixed $u$. We claim that this unique solution is given by 
\begin{align}
    \mathbf{q}_\nu(u,\lambda)=\sum_{j=0}^\infty \mathbf{s}_{j,\nu}(u)\lambda^j,\label{eq:q_nu_series}
\end{align}
where the vector valued functions $\mathbf{s}_{j,\nu}(u),j\ge 0$ are  given by the recursion  
\begin{align}\label{eq:s_nu_def1}
    \mathbf{s}_{0,\nu}(u) &=\frac{n+1}{2}\int_\nu^u e^{(u-s)\mathbf{A}_n} \mathbf{e}_nds+e^{(u-\nu) \mathbf{A}_n }\cdot \mathbf{f}_n,\\
    \mathbf{s}_{k,\nu}(u) &= i \frac{\beta}{4}e^{\frac{\beta}{4} u}\int_\nu^ue^{n/2 (s-u)}e^{(u-s)\mathbf{A}_n}\mathbf{B}_n \mathbf{s}_{k-1,\nu}(s)ds.\label{eq:s_nu_def2}
\end{align}
Lemma \ref{lem:coeff_s} below shows that for $k\ge 0$ the $\CC^n$ valued function $\mathbf{s}_{k,\nu}(u)$ are bounded in norm by a $k$-dependent constant multiple of $e^{\frac{\beta}{4} k u}$, and that these functions satisfy the ODEs
\begin{align}\label{eq:coeff_s_ode1}
     \mathbf{s}_{0,\nu}'(u)&= \mathbf{A}_n  \mathbf{s}_{0,\nu}(u)+ \frac{n+1}{2}\mathbf{e}_n,\\
        \mathbf{s}_{k,\nu}'(u)&= \mathbf{A}_n \mathbf{s}_{k,\nu}(u)+ i \tfrac{\beta}{4}e^{\frac{\beta}{4}u}\mathbf{B}_n \mathbf{s}_{k-1,\nu}(u), \qquad k\ge 1.\label{eq:coeff_s_ode2}
\end{align}
From this it follows that the series given by  \eqref{eq:q_nu_series} can be differentiated in $u$ term-by-term and the resulting series is absolutely convergent. Moreover, the series in \eqref{eq:q_nu_series} satisfies the ODE system \eqref{eq:q_nu_ODE} and hence it is the unique solution of that system.  Lemma \ref{lem:coeff_s} also shows that for any $u>-\infty$ we have $\lim\limits_{\nu\to -\infty} \mathbf{s}_{k,\nu}(u)=e^{k \frac{\beta}{4}u} \mathbf{s}_{k}$ where $\mathbf{s}_{k}$ satisfies \eqref{eq:q_coeff_1}, and that 
\begin{align}\label{eq:q_nu_lim}
    \lim_{\nu\to -\infty} \mathbf{q}_{\nu}(u,\lambda)=\sum_{j=0}^\infty \mathbf{s}_j e^{\frac{\beta}{4} j u} \lambda^j.
\end{align}
From this, we obtain \eqref{eq:q_expansion}, and direct differentiation (which is allowed by Lemma \ref{lem:coeff_s}) gives \eqref{eq:ode_q_vec} as well.
\end{proof}   


\begin{lemma}\label{lem:coeff_s} Consider the functions $ \mathbf{s}_{j,\nu}(u), j\ge 0, u\ge \nu$ defined in \eqref{eq:s_nu_def1}. For every $\nu\in \R$, these are well-defined, and they satisfy the  ODEs given in \eqref{eq:coeff_s_ode1} and   \eqref{eq:coeff_s_ode2}. Moreover, we have the following norm bounds for $\nu\le u$:
\begin{align}\label{eq:coeff_norm_bound}
        \|\mathbf{s}_{k,\nu}(u)\|\le e^{\frac{\beta}{4} k u} \,  \kappa^k \prod_{j=1}^k\frac{n}{j+\frac{4}{\beta}n},
\end{align}
with some $\kappa>0$ (which only depends on $n$). 
Moreover, these functions satisfy 
\begin{align}\label{eq:coeff_nu_lim}
    \lim\limits_{\nu\to -\infty} \mathbf{s}_{k,\nu}(u)=e^{k \frac{\beta}{4}u} \mathbf{s}_{k}.
\end{align}
With $\mathbf{q}_{\nu}(u)$ given by \eqref{eq:q_nu_series} we also have \eqref{eq:q_nu_lim}.
\end{lemma}
We will also need the following lemma.
\begin{lemma}\label{lem:An_eval}
The matrix $\mathbf{A}_n$ has  $n$ distinct  negative eigenvalues given by  $\gamma_{k,n}=\frac{k(k-1)}{2}-\frac{n(n+1)}{2}$, $1\le k\le n$. 
\end{lemma}
\begin{proof}
 Our starting point is the following version of the Chu-Vandermonde identity that holds for non-negative integers $a, b, r$ (see e.g.~(5.24) in \cite{Concrete}
).
 \begin{align}\label{eq:chu-v}
     \sum_{k=0}^a (-1)^k \binom{a}{k}\binom{k+r}{b}=(-1)^a \binom{r}{b-r}.
 \end{align}
Here $\binom{x}{y}=0$ if $y<0$ or $x<y$. Now consider the $n\times n$ matrix $\mathbf{T}_n$ defined as $[\mathbf{T}_n]_{a,b}=\binom{a}{b}(-1)^b$, note that this is a lower triangular matrix. Identity \eqref{eq:chu-v} with $r=0$ shows that $\mathbf{T}_n=\mathbf{T}_n^{-1}$. We claim that $\mathbf{T}_n \mathbf{A}_n \mathbf{T}_n$ is an upper bidiagonal matrix with entries
\begin{align}\label{eq:TAT}
 [\mathbf{T}_n \mathbf{A}_n \mathbf{T}_n]_{k,k}= \frac{k(k-1)}{2}-k n, \qquad [\mathbf{T}_n \mathbf{A}_n \mathbf{T}_n]_{k,k+1}=\frac{k(n-k)}{2}
\end{align}
 in the $k$th row. Since  $  \gamma_{n+1-k,n}=\frac{k(k-1)}{2}-k n$, from this the statement of the lemma follows.\\
To prove \eqref{eq:TAT} we start with the expression for the  $(\ell,m)$ entry of $\mathbf{T}_n\mathbf{A}_n\mathbf{T}_n$:
\begin{align}\label{eq:TAT1}
    [\mathbf{T}_n\mathbf{A}_n\mathbf{T}_n]_{\ell,m}=\sum_{j=1}^n t_{\ell,j} (a_{j,j} t_{j,m}+a_{j,j-1}t_{j-1,m}+a_{j,j+1}t_{j+1,m}). 
\end{align}
Here $t_{\ell,m}, a_{\ell,m}$ are the respective entries of $\mathbf{T}_n$ and $\mathbf{A}_n$, and they are defined to be zero if we do not have $1\le \ell, m\le n$. 
Using simple identities of binomial coefficients, we can rewrite each term in \eqref{eq:TAT1} as linear combinations of expressions of the form
$(-1)^j\binom{\ell}{j} \binom{j+r}{m}$ with finitely many choices for $r$. By applying  \eqref{eq:chu-v} to each of these terms, a somewhat lengthy but straightforward calculation shows that $\mathbf{T}_n\mathbf{A}_n\mathbf{T}_n$ is indeed upper bidiagonal with the claimed entries. 
\end{proof}




\begin{proof}[Proof of Lemma \ref{lem:coeff_s}]
The ODEs \eqref{eq:coeff_s_ode1} and   \eqref{eq:coeff_s_ode2} can be checked directly by differentiation from \eqref{eq:s_nu_def1} and \eqref{eq:s_nu_def2}.
 From Lemma \ref{lem:An_eval} it follows that the spectral radius of $e^{x \mathbf{A}_n}$ with $x\ge 0$ is $e^{-n x}$. This implies that there is a finite constant  $\kappa>0$ (depending only on $n$) so that
 \[
 \| e^{x \mathbf{A}_n} \mathbf{v}\|\le \kappa e^{-n x} \|\mathbf{v}\|.
 \]
The spectral radius of $\mathbf{B}_n$ is $n$, hence we get $\|\mathbf{B}_n \mathbf{v}\|\le n  \|\mathbf{v}\|$.
From \eqref{eq:s_nu_def1} and \eqref{eq:s_nu_def2}  we get
\begin{align}\label{eq:s_norm_bnd0}
    \|\mathbf{s}_{0,\nu}(u)\|&\le \kappa \frac{n+1}{2n}+\kappa e^{-(u-\nu)n} \|\mathbf{f}_n\|,\\
    \|\mathbf{s}_{k,\nu}(u)\|&\le \kappa n \frac{\beta}{4} e^{\frac{\beta}{4}u} \int_{\nu}^u e^{(s-u)(\frac{\beta}{4}+n)} \|\mathbf{s}_{k-1,\nu}(s)\| ds.
\end{align}
From \eqref{eq:s_norm_bnd0} it follows that $\|\mathbf{s}_{0,\nu}(u)\|$ is bounded by a constant uniformly in $\nu$, and by induction, it follows that there is a finite $\kappa_0>0$ so that 
\begin{align}\label{eq:s_nu_norm_bnd}
    \|\mathbf{s}_{k,\nu}(u)\|\le e^{\frac{\beta}{4} k u} \, \kappa_0 \kappa^k \prod_{j=1}^k\frac{n}{j+\frac{4}{\beta}n} .
\end{align}
Increasing the value of $\kappa$ gives the bound \eqref{eq:coeff_norm_bound}. 
Taking the limit $\nu\to -\infty$ in \eqref{eq:s_nu_def1} gives
\begin{align*}
  \lim_{\nu\to -\infty}   \mathbf{s}_{0,\nu}(u)&=\frac{n+1}{2}\int_{-\infty}^u e^{(u-s)\mathbf{A}_n} \mathbf{e}_nds+\lim_{\nu\to -\infty} e^{(u-\nu) \mathbf{A}_n }\mathbf{f}_n\\
  &=\frac{n+1}{2}\int_{-\infty}^0 e^{t \mathbf{A}_n} \mathbf{e}_ndt=-\frac{n+1}{2} \mathbf{A}_n^{-1} \mathbf{e}_n,
\end{align*}
proving the $k=0$ case of \eqref{eq:coeff_nu_lim}. 
Note that  
 $\mathbf{A}_n^{-1}$ is well-defined by Lemma \ref{lem:An_eval}. We can now prove  \eqref{eq:coeff_nu_lim} for general $k\ge 1$  by induction. Assume that \eqref{eq:coeff_nu_lim} holds for $k-1$ for some $k\ge 1$. 
By \eqref{eq:s_nu_norm_bnd} we can use dominated convergence to get
\begin{align*}
  \lim_{\nu\to -\infty}   \mathbf{s}_{k,\nu}(u)&=i \frac{\beta}{4}e^{\frac{\beta}{4} u}\int_{-\infty}^ue^{\beta/4 (s-u)}e^{(u-s)\mathbf{A}_n}\mathbf{B}_n e^{(k-1) \frac{\beta}{4}s}\mathbf{s}_{k-1}ds\\
  &=i \frac{\beta}{4}e^{\frac{k \beta}{4} u}\int_{-\infty}^ue^{k \frac{\beta}{4} (s-u)}e^{(u-s)\mathbf{A}_n}\mathbf{B}_n \mathbf{s}_{k-1}ds=i \frac{\beta}{4}e^{\frac{k \beta}{4} u}\int_{-\infty}^0e^{(-k \frac{\beta}{4}+\mathbf{A}_n) t }dt \, \mathbf{B}_n\mathbf{s}_{k-1}
\\&=i e^{\frac{k \beta}{4} u}\left(k-\tfrac{4}{\beta}\mathbf{A}_n\right)^{-1} \, \mathbf{B}_n\mathbf{s}_{k-1}.
\end{align*}
Again, $\left(k \mathbf{I}-\tfrac{4}{\beta}\mathbf{A}_n\right)^{-1}$ is well-defined because of Lemma \ref{lem:An_eval}. This concludes the proof of \eqref{eq:coeff_nu_lim}. Finally, we use the norm bound \eqref{eq:coeff_norm_bound} again to conclude that 
\[
 \lim_{\nu\to -\infty} \mathbf{q}_{\nu}(u,\lambda)=
\sum_{j=0}^\infty \lim_{\nu\to -\infty} \mathbf{s}_{j,\nu}(u) \lambda^j
 =
 \sum_{j=0}^\infty \mathbf{s}_j e^{\frac{\beta}{4} j u} \lambda^j,
\]
proving \eqref{eq:q_nu_lim}.
\end{proof}
To obtain Corollary \ref{cor:int_case_rho} and Proposition \ref{prop:small_lambda_asym}, we prove the following claim. 

\begin{claim} We have the following identities:
   \label{cl:vAnBnf}
% \yahui{$\quad \mathbf{v}\cdot \mathbf{A}_n = -n \mathbf{A}_n$  we missed this - at beginning}
    \begin{align}
      \label{eq:AnBnf}
            &\mathbf{s}_0= \mathbf{f}_n, \\
            &\mathbf{A}_n \mathbf{B}^{\ell}_n \mathbf{f}_n=\sum_{j=0}^{\lfloor \frac{\ell-1}{2}\rfloor} \left(\tbinom{\ell}{2j+2}-n\tbinom{\ell}{2j+1}\right)\mathbf{B}^{\ell-2j}_n\mathbf{f}_n, \quad \text{ for all $1\le \ell \le 2n$}
           \label{eq:AnBnf_2}
          \\
          \label{eq:vBnf}
          &\mathbf{v}_n\cdot \mathbf{B}^{2j}_n \mathbf{f}_n = 0,\;\; \text{ for all $1\le j<n$, \quad and }\qquad \mathbf{v}_n\cdot \mathbf{B}^{2n}_n \mathbf{f}_n = (-1)^n \binom{2n}{n}^{-1}\frac{(2n)!}{2}.
    \end{align}
    \end{claim}
\begin{proof}
From \eqref{eq:AB} and \eqref{eq:f} we have, 
\begin{align}\label{eq:AB1}
[\mathbf{B}^{\ell}_n \mathbf{f}_n]_k &= k^{\ell}, \quad
[\mathbf{A}_n \mathbf{f}_n]_k =-\frac{n+1}{2} \ind(k=1),
\qquad 1\le k\le n,\\
    [\mathbf{A}_n \mathbf{B}^\ell_n\mathbf{f}_n]_k&=\frac{k(k+n)}{2}(k-1)^\ell-k^{2+\ell}+\frac{k(k-n)}{2}(k+1)^\ell, \qquad 1\le k\le n.\label{eq:AB2}
\end{align}
%Note that  last equation works for $k=1$ and $k=n$ as well. 
From \eqref{eq:AB1} we have $\mathbf{A}_n \mathbf{f}_n=-\tfrac{n+1}{2}\mathbf{e}_n$, hence from \eqref{eq:q_coeff_1} we obtain  \eqref{eq:AnBnf}. Expanding $(k\pm 1)^\ell$ and checking the coefficients of $k$ yields \eqref{eq:AnBnf_2}. \\
To prove \eqref{eq:vBnf}, we start with the identity
    \begin{align*}
        \mathbf{v}_n\cdot \mathbf{B}^{2 j}_n \mathbf{f}_n  = \tbinom{2n}{n}^{-1} \sum_{k=1}^n (-1)^k \tbinom{2n}{n+k} k^{2j}  = \frac12 (-1)^n\tbinom{2n}{n}^{-1} \sum_{k=0}^{2n}(-1)^{2n-k} \binom{2n}{k} (k-n)^{2j}.
    \end{align*}
The sum on the right is just the result of the finite difference operator \[\Delta f(x)=f(x+1)-f(x)\] applied $2n$ times to the function  $f(x)=(x-n)^{2j}$, and then evaluated at $x=0$. 
Applying $\Delta^{2n}$ to the polynomial $x^{2n}$ yields the constant $(2n)!$, and applying the same operator to any polynomial with degree less than $2n$ yields the constant 0.  From this, we obtain \eqref{eq:vBnf}.
\end{proof}





\begin{proof}[Proof of Corollary \ref{cor:int_case_rho}]
The equation 
\[
\rho^{(1)}_{\beta,n}(\lambda) =\frac{1}{2\pi}\left(1+2\;\mathbf{v}_n^T\Re \mathbf{q}(\lambda)\right)
\]
follows from \eqref{eq:AnBnf} and Theorem \ref{thm:even_beta}. To prove the second part of \eqref{eq:int_case_rho}, we note that from \eqref{cl:vAnBnf} we have $\mathbf{s}_0 = \mathbf{f}_n$. From $(1-1)^{2n} = 0$ we also get 
\[
\mathbf{v}_n \cdot \mathbf{s}_0 = \mathbf{v}_n \cdot \mathbf{f}_n = \sum_{k=1}^n (-1)^k \frac{\binom{2n}{n+k}}{\binom{2n}{n}} = -\frac{1}{2}.
\]
From the definition of $\mathbf{s}_j$ it follows that  $\mathbf{s}_j$ is pure imaginary when j is odd, and real when j is even. Hence we get
\begin{align*}
      &\rho^{(1)}_{\beta,n}(\lambda)  =\frac{1}{2\pi}\left(1+2\;\mathbf{v}_n^T\Re \mathbf{q}(\lambda)\right)=\frac{1}{2\pi}(1+ 2\sum_{j=0}^\infty \mathbf{v}_n\cdot \mathbf{s}_{2j} \lambda^{2j}) =\frac{1}{\pi} \sum_{j=1}^\infty \mathbf{v}_n\cdot \mathbf{s}_{2j} \lambda^{2j},
\end{align*}
which completes the proof. 
\end{proof}

\begin{proof}[Proof of Proposition \ref{prop:small_lambda_asym}]
    Following \eqref{eq:int_case_rho}, it suffices to show that
     \begin{align}
         \mathbf{v}\cdot\mathbf{s}_{2\ell}=0,\quad \text{for $1\le \ell\le n-1$},\quad\text{and} \quad 
         \mathbf{v}\cdot\mathbf{s}_{2n}=   \frac12 \binom{2n}{n}^{-1}\frac{\left(\tfrac{\beta}{2}\right)^{2n}}{\left(1+\frac{\beta}{2}\right)^{\uparrow 2n}}.
     \end{align}
Fix $1\le \ell\le 2n$. Consider the vector space $\mathcal{V}_\ell$ spanned by the vectors $\mathbf{B}_n^{\ell-2j}\mathbf{f}_n$, with $j=0,\dots,  \lfloor\frac{\ell-1}{2}\rfloor$.
Equation \eqref{eq:AnBnf_2} shows that $\mathbf{A}_n$ maps $\mathcal{V}_\ell$ to itself and acts as a lower-triangular matrix with respect to this basis. This implies that for $\gamma>0$ the mapping $\mathbf{I}-\gamma \mathbf{A}_n: \mathcal{V}_\ell \mapsto \mathcal{V}_\ell$ is invertible, and the inverse is a lower triangular matrix with respect to our basis. 
In other words, there exist constants $ c_{\ell,\gamma, m, n} $ such that,
\begin{align}
    \label{eq:invBnf}
    (\mathbf{I}-\gamma \mathbf{A}_n)^{-1}   \mathbf{B}^{\ell}_n   \mathbf{f}_n = \sum_{m=0}^{\lfloor \frac{\ell-1}{2} \rfloor} c_{\ell,\gamma, m, n}  \mathbf{B}^{\ell-2m}_n \mathbf{f}_n.
\end{align}
Multiplying both sides by $(\mathbf{I}-\gamma \mathbf{A}_n)$ and matching the corresponding coefficients of $\mathbf{B}_n^{\ell} \mathbf{f}_n$ term, we get
\begin{align} \label{eq:invBnf_coeff}
 c_{\ell,\gamma, 0, n}=(1+ \gamma \ell (n-\tfrac{\ell-1}{2}))^{-1}.   
\end{align}
 By definition, we have for all $1\leq \ell \leq n,$ 
\begin{align*}
    \mathbf{s}_{2\ell} &= (-1)^{\ell} 
(2\ell \,\mathbf{I}-\tfrac{4}{\beta} \mathbf{A}_n)^{-1}\mathbf{B}_n ((2\ell-1) \mathbf{I}-\tfrac{4}{\beta} \mathbf{A}_n)^{-1}\mathbf{B}_n\cdots ( \mathbf{I}-\tfrac{4}{\beta} \mathbf{A}_n)^{-1}\mathbf{B}_n \mathbf{f}_n\\
 &= (-1)^{\ell} \frac{1}{(2\ell)!}
( \mathbf{I}-\tfrac{4}{\beta(2\ell)} \mathbf{A}_n)^{-1}\mathbf{B}_n (\mathbf{I}-\tfrac{4}{\beta (2\ell-1)} \mathbf{A}_n)^{-1}\mathbf{B}_n\cdots ( \mathbf{I}-\tfrac{4}{\beta} \mathbf{A}_n)^{-1}\mathbf{B}_n \mathbf{f}_n.
   %  \left((k-\frac{4}{\beta} \mathbf{A}_n)^{-1}\mathbf{B}_n \right)\mathbf{f}_n\\
   % &=(-1)^\ell \frac{1}{(2\ell)!} \mathbf{v}_n\cdot\prod_{k=1}^{2\ell}\left((\mathbf{I}-\frac{4}{\beta k} \mathbf{A}_n)^{-1}\mathbf{B}_n \right)\mathbf{f}_n
   \end{align*}
%    \\
%    \nonumber\\& \qquad=(-1)^\ell \prod_{k=1}^{2j} \frac{1}{k}c_{k,\frac{4}{\beta k },0,n}  \mathbf{v}_n \cdot \mathbf{B}_n^{2j} f + \sum_{k=1}^{l-1} \tilde{c}_{2j,\frac{4}{\beta},k,n} \mathbf{v}_n\cdot \mathbf{B}_n^{2j-2k} f,
% \end{align}
Applying  \eqref{eq:invBnf} repeatedly (from the right), we obtain that 
\[
  \mathbf{s}_{2\ell} =(-1)^\ell \frac{1}{(2\ell)!}\prod_{k=1}^{2\ell} c_{k,\frac{4}{\beta k },0,n} \,\mathbf{B}_n^{2\ell} \mathbf{f}_n + \sum_{k=1}^{\ell-1} \tilde{c}_{\ell, k, n} \, \mathbf{B}_n^{2\ell-2k} \mathbf{f}_n,
\]
with some finite coefficients  $\tilde{c}_{\ell, k, n}$.
Finally, by \eqref{eq:vBnf} and \eqref{eq:invBnf_coeff}, we conclude that
\begin{align*}
     \mathbf{v}_n\cdot \mathbf{s}_{2\ell} = 0 \quad \text{for } 1\le \ell\le n-1,
\end{align*}
and
\begin{align*}
    \mathbf{v}_n\cdot \mathbf{s}_{2n}
    = \frac{1}{2}\binom{2n}{n}^{-1} \prod_{j=1}^{2n}\left(1+\tfrac{4}{\beta}\Big(n-\tfrac{j-1}{2}\Big)\right)^{-1}.
\end{align*}
%=\frac{\beta^{2n}\Gamma(n+1)^2\Gamma\left(1+\frac{\beta}{2}\right)}
The last constant can also be written as 
\[
\frac12 \binom{2n}{n}^{-1}\frac{(\tfrac{\beta}{2})^{2n}}{\left(1+\frac{\beta}{2}\right)^{\uparrow 2n}},
\]
which completes the proof of Proposition \ref{prop:small_lambda_asym}.
\end{proof}

\section{Additional results}\label{sec:additional}


\subsection{ODE system for general $\beta>0$}

Theorem \ref{thm:even_beta} has the following extension. 

\begin{proposition}\label{prop:ODE_gen}
Fix $\beta>0, \delta>0$. Let $q_k(\lambda)=E[e^{i k \alpha_{\lambda}(0)}]$ where $\alpha_\lambda(u)$ is the strong solution of \eqref{eq:alphaSDE} with initial condition \eqref{eq:SDE_initial}. Then $q_k(\lambda)$ is differentiable for $\lambda\in \R$, $\lambda\neq 0$, and we have
\begin{align}\label{eq:genODE1}
     \tfrac{\beta}{4}\lambda q_{1}'(\lambda)&=\tfrac{1}{2}(1+\delta)+(i \lambda \tfrac{\beta}{4}-1) q_{1}(\lambda)+\tfrac{1}{2}(1-\delta)q_{2}(\lambda),\\\label{eq:genODE2}
   \tfrac{\beta}{4} \lambda q'_{k}(\lambda)&=\tfrac{1}{2} k(k+\delta) q_{k-1}(\lambda)+(i k \lambda \tfrac{\beta}{4}-k^2)q_{k}(\lambda)+\tfrac{1}{2} k (k-\delta )q_{k+1}(\lambda), \qquad 2\le k,\\
    \lim_{\lambda\to 0}q_{k}(\lambda)&=1. \label{eq:genODE3}
\end{align}    
\end{proposition}
\begin{proof}
Let $\tilde q_k(u,\lambda)=E[e^{i k \alpha_\lambda(u)}]$ for $k\ge 1$. 
As argued before, we can take expectations in \eqref{eq:QkSDE} to get differential equations for $\tilde q_k(u,\lambda)$ in $u$ for $u\le 0$. By the scale invariance property \eqref{eq:scale_inv} we have 
\[
\tilde q_k(u,\lambda)=q_k(e^{\frac{\beta}{4} u} \lambda),
\]
which allows us to rewrite the differential equations for $\tilde q_k(u,\lambda)$, $\lambda\neq 0$ 
in terms of $q_k(\lambda)$, yielding the differential equations  \eqref{eq:genODE1} and \eqref{eq:genODE2}. The boundary condition \eqref{eq:genODE3} follows from the bounded convergence theorem and the continuity of $\alpha_\lambda$ in $\lambda$. 
\end{proof}




    
\subsection{Explicit solutions}

We have already shown how we can obtain the Sine-kernel, the pair correlation of the $\Sine_2$ process from Theorem \ref{thm:sineb_series}. In this section, we demonstrate how we can recover the known pair correlation function \eqref{eq:beta=4} for the $\operatorname{Sine}_4$ process, and also provide an explicit formula for the density of the $\HPb$ process for $\delta=1$. 




When $n=2$ from \eqref{eq:AB} and \eqref{eq:vectors} we get 
\[
\mathbf{A}_2=\mat{-1}{-\frac12}{4}{-4}, \quad \mathbf{B}_2=\mat{1}{0}{0}{2}, \quad \mathbf{v}_2=[-\tfrac23, \tfrac16]^T.
\]
Hence, by Theorem \ref{thm:sineb_series} and Theorem \ref{thm:even_beta}, we have the following representation for the pair correlation of the  $\Sine_4$ process: 
\begin{align}\label{eq:rho4}
     \rho^{(2)}_{4}(0,\lambda)=\frac{1}{4\pi^2}+\frac1{2\pi^2}\left(-\frac{2}{3} \Re[q_1(\lambda)]+\frac16\Re[q_2(\lambda)]\right),
\end{align}
where $q_1, q_2$ solves the ODE system 
\begin{align}\label{eq:ODE4_1}
    \lambda q_1'(\lambda)&=\frac32-q_1(\lambda)-\frac12 q_2(\lambda)+i \lambda q_1(\lambda),\qquad q_1(0)=1,\\
    \lambda q_2'(\lambda)&=4 q_1(\lambda)-4 q_2(\lambda)+2 i \lambda q_2(\lambda), \qquad \phantom{....}q_2(0)=1.\label{eq:ODE4_2}
\end{align}
This system can be reduced to the following second-order ODE for $q_2$:
\begin{align*}
    % q_1(\lambda)&=\frac{1}{4} \lambda  q_2'(\lambda )+\frac{1}{2} (2-i \lambda ) q_2(\lambda ),\\
    \lambda^2 q_2''(\lambda)&=\left(-6 \lambda +3 i \lambda ^2\right) q_2'(\lambda )+\left(2 \lambda ^2+8 i \lambda -6\right) q_2(\lambda )+6,\qquad q_2(0)=1, q_2'(0)=\frac{2 i}{3},
\end{align*}
which  has the unique solution
% With $q_2(\lambda)=e^{i \lambda} \lambda^{-2} r_2(\lambda)$ we get
% \[
% 6 \lambda -e^{i \lambda } \lambda r_2''(\lambda )+\left(i e^{i \lambda } \lambda -2 e^{i \lambda }\right) r_2'(\lambda )=0.
% \]
% This can be solved (noting that $r_2'(0)=0$) to get
% \[
% r_2'(\lambda)=-\frac{3 \sin (\lambda )}{\lambda ^2}+3 i e^{-i \lambda }+\frac{3 e^{-i \lambda }}{\lambda }.
% \]
% From this, we can recover $q_2$ to get 
% \[
% q_2(\lambda)=e^{i \lambda} \lambda^{-2}\int_0^\lambda \left(\frac{3 \cos (x)}{x}-\frac{3 \sin (x)}{x^2}+3 \sin (x)-\frac{3 i \sin (x)}{x}+3 i \cos (x) \right)dx.
% \]
% OR
\[
q_2(\lambda)=
%\tfrac32 i e^{i \lambda} \lambda^{-2} \int_0^\lambda\frac{ e^{-i x} \left(e^{2 i x}-1-2 i x+ 2x^2\right)}{ x^2}dx
\frac{3i(1+2i\lambda-e^{2i \lambda})}{2\lambda^3}-3i e^{i \lambda}\lambda^{-2}\int_0^\lambda \frac{\sin(x)}{x}dx.
\]
We can now obtain $q_1(\lambda)$ from \eqref{eq:ODE4_2}, 
and substituting these into  \eqref{eq:rho4} yields
the classical formula  \eqref{eq:beta=4} for  $\rho^{(2)}_{4}(0,\lambda)$.
\medskip

To obtain the density of the $\operatorname{HP}_{\beta,1}$ process we can use Theorem \ref{thm:even_beta} or Corollary \ref{cor:int_case_rho}. Here $n=1$, and we have $\mathbf{A}_1=-1$, $\mathbf{B}_1=1$, $\mathbf{v}_1=-\tfrac{1}{2}$.  Hence  $\mathbf{s}_k$ satisfies the scalar recursion 
\[
\mathbf{s}_0=1, \qquad \mathbf{s}_k=i (k+\tfrac{4}{\beta})^{-1} \mathbf{s}_{k-1}, \qquad k\ge 1
\]
which yields $\mathbf{s}_k=\frac{i^k}{\left(1+\tfrac{4}{\beta}\right)^{\uparrow k}}$.  Corollary \ref{cor:int_case_rho} gives
\begin{align}
   \rho_{\beta,1}^{(1)}(\lambda)=-\frac12 \sum_{k=1}^\infty \frac{(-1)^k }{\left(1+\tfrac{4}{\beta}\right)^{\uparrow 2k}}
   \lambda^{2k}
   =
   \frac{1 }{4 \pi (1+ \frac{2}{\beta}) (1+\frac{4}{\beta})} \, _1F_2\left(1;\frac{3}{2}+\frac{2}{\beta },2+\frac{2}{\beta };-\frac{\lambda ^2}{4}\right) \lambda ^2.
\end{align}
This expression can also be obtained by solving the ODE given by  Theorem \ref{thm:even_beta}:
\begin{align*}
    \tfrac{\beta}{4}\lambda q'(\lambda) &= (i\tfrac{\beta}{4} \lambda -1) q(\lambda) + 1 , \quad q(0) = 1.\\
    q(\lambda)
&= \frac{4}{\beta}\,e^{i\lambda}\,\lambda^{-4/\beta}
\int_{0}^{\lambda} s^{\,4/\beta - 1}\,e^{-i s}\,ds.
\end{align*}
Using the first equation in Corollary \ref{cor:int_case_rho} now yields the equivalent expression
\begin{align*}
     \rho_{\beta,1}^{(1)}(\lambda) = \frac{1}{2\pi}-\frac{1}{2\pi}\Re{q(\lambda)}=
\frac{1}{2\pi}-\frac{2}{\beta \pi}\lambda^{-4/\beta} \int_0^\lambda s^{4/\beta-1} \cos(\lambda-s) ds.
\end{align*}



\subsection{ODE representation for $\rho_{2n}^{(2)}(0,\lambda)$}

Theorem \ref{thm:even_beta} and Corollary \ref{cor:int_case_rho} with $\beta=2n$ provides a description for $\rho_{2n}^{(2)}(0,\lambda)$ as 
\[
\rho_{2n}^{(2)}(0,\lambda)=\frac{1}{4\pi^2}\left(1+2\;\mathbf{v}_n^T\Re \mathbf{q}(\lambda)\right)
\]
with $\mathbf{q}$ as the solution of the ODE \eqref{eq:ode_q_vec}. Motivated by \cite{Forrester_20212}, one can ask if there is a way to represent $\rho_{2n}^{(2)}(0,\lambda)$ itself as the solution of a differential equation using our representation. It turns out that the function $r_n(\lambda)=2 \mathbf{v}_n^T \mathbf{q}(\lambda)$ solves an order $n$ complex-valued linear differential equation where the coefficients are polynomials with degree at most $n$. This gives a representation for  $\rho_{2n}^{(2)}(0,\lambda)$ as the \emph{real part} of the solution of an order $n$  linear differential equation. 

We will only give the outline of the proof of this statement, without spelling out the computational details. 
It can be shown that $\mathbf{v}_n^T$ is the left eigenvector of $\mathbf{A}_n$ corresponding to the eigenvalue $-n$. From Lemma \ref{lem:An_eval} it follows that  $\mathbf{A}_n$ has $n$ linearly independent left eigenvectors, let $\mathbf{w}_k^T, 1\le k\le n-1$ denote the  ones besides $\mathbf{v}_n^T$. Introducing the new variables $r_k(\lambda)=\mathbf{w}_k^T \mathbf{q}(\lambda)$, $1\le k\le n-1$, we see that $\mathbf{r}(\lambda)=[r_1(\lambda),\dots,r_n(\lambda)]^T$ satisfies a transformed version of \eqref{eq:ode_q_vec}, with $\mathbf{B}_n, \mathbf{A}_n, \mathbf{e}_n, \mathbf{f}_n$ replaced with their respective representations in a new basis. In this new representation $\mathbf{A}_n$ will become a diagonal matrix, and by choosing the normalization in the eigenvectors $\mathbf{w}_k, 1\le k\le n-1$ appropriately, one can show that $\mathbf{B}_n$ will turn into a tridiagonal matrix with explicitly computable entries. This allows us to successively eliminate $r_{n-1}, r_{n-2}, \dots, r_1$ from the ODE system for $\mathbf{r}(\lambda)$, leading to an order $n$ linear ODE  with polynomial coefficients of degree at most $n$. The initial condition for $\mathbf{r}(0)$ will be transformed into equations for $r_n(0), r'_n(0), \dots, r_n^{(n-1)}(0)$. 


A similar (but computationally more involved) approach could be taken to obtain a differential equation for which $\rho_{2n}^{(2)}(0,\lambda)$ itself is a solution.  Consider the power series representation of $\rho_{2n}^{(2)}(0,\lambda)$ from Theorem \ref{thm:sineb_series}. Let $\mathbf{z}(\lambda)=\sum_{j=1}^\infty \mathbf{s}_{2j} \lambda^{2j}$. Using the recursive definition of $\mathbf{s}_{k}, k\ge 0$ we can set up a second order vector valued linear differential equation for $\mathbf{z}(\lambda)$ with terms of the form $\lambda^2 \mathbf{z}''(\lambda)$, $\lambda \mathbf{z}'(\lambda)$, $\lambda \mathbf{z}(\lambda)$, and $\mathbf{z}(\lambda)$ (with matrix valued coefficients). Using the arguments of the previous paragraph it might be possible to turn this vector valued second order ODE into an order $2n$ linear ODE  for  $ \rho_{2n}^{(2)}(0,\lambda)=\frac{1}{2\pi^2} \mathbf{v}_n^T \cdot \mathbf{z}(\lambda)$. However, we only carried out this reduction for small values of $n$ with the assistance of computer algebra. 


\subsection{Regularity of the pair correlation function $\rho_{\beta}^{(2)}(0,\lambda)$}

We turn to the proof of Theorem \ref{thm:beta_cont}. We will use the following proposition, which can be shown using standard results on SDE flows (see e.g.~Theorem 49 in \cite{Protter}).

\begin{proposition}\label{prop:sde_cont}
Let $Z$ be a two-sided Brownian motion.
There is a process $\alpha_{\lambda,\beta}(u)$, $\lambda\in \R$, $\beta>0$, $u\in \R$ that is adapted to the filtration of $Z$ as a process in $u$, analytic in $\lambda$ for fixed $u$ and  $\beta$, continuous in  $\beta$ for fixed $u$ and $\lambda$, and has the same distribution as the unique strong solution $u\mapsto \alpha_{\lambda}(u)$ of  \eqref{eq:alphaSDE_b} for a fixed $\beta$.     
\end{proposition}

\begin{proof}[Proof of Theorem \ref{thm:beta_cont}]
Consider the process $\alpha_{\lambda,\beta}(u)$ given by  Proposition \ref{prop:sde_cont}, and set
\[
q_k(\beta,\lambda)=E [\exp(i k \alpha_{\lambda,\beta}(0))].
\]
From Proposition \ref{prop:sde_cont} it follows that these functions are analytic in $\lambda$ and continuous in $\beta$, and by definition we have $|q_k|\le 1$.  By Theorem \ref{thm:sineb_corr} we have
\begin{align}\label{eq:qqq}
    \rho_\beta^{(2)}(0,\lambda)=\frac{1}{4\pi^2}+\frac{1}{2\pi^2}\sum_{k=1}^\infty  \frac{(-\beta/2)^{\uparrow k}}{(1+\beta/2)^{\uparrow k}} \Re q_k(\beta,\lambda).
\end{align}
The first part of the theorem now follows:  the partial sums in \eqref{eq:qqq} are continuous in $\lambda$ and in $\beta$, and these continuous functions converge uniformly to  $\rho_\beta^{(2)}(0,\lambda)$ because of the summable bound \eqref{eq:coeff_bnd}.


To see that for $\beta> 2, \lambda>0$ we have joint continuity in $\lambda$ and $\beta$ we use the following statement: if $f(x,y)$ is continuous in both variables, and it is also locally uniformly Lipschitz in $y$ then $f(x,y)$ is jointly continuous. Indeed, if $(x_n,y_n)\to (x,y)$ then 
\[
|f(x,y)-f(x_n, y_n)|\le |f(x,y)-f(x_n,y)|+|f(x_n,y)-f(x_n,y_n)|\to 0,
\]
where we use the continuity in $x$ in the first term and the uniform Lipschitz bound in the 
second term. 


From Proposition \ref{prop:ODE_gen} we see that $q_k(\beta,\lambda)$ satisfies the ODEs \eqref{eq:genODE1} and \eqref{eq:genODE1} with $\delta=\frac{\beta}{2}$. From this, it follows that for $k\ge 1$, $\lambda>0$ we have
\[
|\partial_\lambda \Re q_k(\beta,\lambda)|\le c(k^2 \lambda^{-1}+k)
\]
with a $\beta$-dependent constant $c$, that is locally bounded in $\beta$. Together with the bound \eqref{eq:coeff_bnd} this implies that for $\beta>2$, $\lambda>0$ the function $ \rho_\beta^{(2)}(0,\lambda)$ is locally uniformly Lipschitz in $\lambda$, completing the proof. 
% By Proposition \ref{prop:sde_cont}
%  we can realize a solution $\alpha: (-\infty,0]\times \R_{+}\times \R\mapsto \R$ of \eqref{eq:alphaSDE} with $\delta=\beta/2$ that is continuous in $\beta, \lambda$ and $t$. Then, for a fixed $\lambda$, $k>0$, the function 
% \[
% \beta\mapsto E_{\beta,\beta/2} [\cos k \alpha_\lambda(0)]
% \]
% is continuous. 
% \yahui{Since the summation in \eqref{eq:twopoint_cor} is bounded, and each summand is continuous in $\beta,\lambda$}
% From this, the statement follows.\benedek{add more detail?} 
\end{proof}


\section{Open problems} \label{sec:open}

We finish with a list of open problems related to our results. \medskip


Our estimates in the proof of   Theorem \ref{thm:beta_cont} are far from optimal. Potentially,  a more careful asymptotic analysis of $E[e^{i k \alpha_\lambda(0)}]$ could produce bounds that match the conjectured order $\lambda^{-\min(2,4/\beta)}$  for the truncated pair correlation.  
\begin{problem}
  Improve the bounds in Theorem \ref{thm:beta_cont} for the truncated pair correlation function of the $\Sineb$ process for general $\beta>0$.
\end{problem}


Proposition \ref{prop:ODE_gen} shows that $q_k(\lambda)=E[e^{i k \alpha_\lambda(0)}]$ solves an infinite system of ODEs, but it is not clear that the solution is unique. A uniqueness result would provide a self-contained characterization of the pair correlation function $\rho_{\beta}^{(2)}(0,\lambda)$. 


\begin{problem}
Show that the infinite ODE system \eqref{eq:genODE1}-\eqref{eq:genODE3}  (with possibly some additional constraints) has a unique solution.
\end{problem}

Our proof of Theorem \ref{thm:1p_HP} relies on $\delta$ being a positive real. A natural question to ask if the result could be extended to more general $\delta$ values.
\begin{problem}
    Extend Theorem \ref{thm:1p_HP} for the density of the $\HPb$ process for complex $\delta$ with $\Re \delta>-1/2$.
\end{problem}

It would also be interesting to extend the second part of Theorem \ref{thm:beta_cont} for all $\beta>0$.
\begin{problem}
    Prove (or disprove) that $\rho_{\beta}^{(2)}(0,\lambda)$ is jointly continuous in $\lambda$ and $\beta$ for $\lambda\in \R$, $\beta>0$. 
\end{problem}

The small $\lambda$ asymptotics of Corollary \ref{cor:small_lambda} should hold for all $\beta>0$, but we did not find a reference for this in the literature.  

\begin{problem}
Extend the small $\lambda$ asymptotics of Corollary \ref{cor:small_lambda} for all $\beta>0$ with the factorials replaced with Gamma functions in \eqref{eq:small_lambda}.
\end{problem}

As mentioned before, it would be instructive to show directly that the representation for $\rho^{(2)}_{2n}(0,\lambda)$ found in Theorem \ref{thm:sineb_series} agrees with \eqref{eq:paircorr_For}.

\begin{problem}
    Show directly that the power series given in Theorem \ref{thm:sineb_series} is equivalent to the integral expression \eqref{eq:paircorr_For} proved by Forrester.
\end{problem}

Finally, the most challenging problem is to extend the results of this paper to provide any information about higher-order correlation functions.


\begin{problem}
Find a characterization of the higher order correlation functions of the $\Sineb$ process in terms of the $\alpha_\lambda$ diffusion given in \eqref{eq:alphaSDE_b}.
\end{problem}






%\bibliography{rmt}


\def\cprime{$'$}
\begin{thebibliography}{10}

\bibitem{AGZ}
G.~Anderson, A.~Guionnet, and O.~Zeitouni.
\newblock {\em Introduction to random matrices}.
\newblock Cambridge University Press, 2009.

\bibitem{andrews1999special}
G.E. Andrews, R.~Askey, and R.~Roy.
\newblock {\em Special {F}unctions}.
\newblock Encyclopedia of Mathematics and its Applications. Cambridge
  University Press, 1999.

\bibitem{Spike1}
A.~Bloemendal and B.~Vir{\'a}g.
\newblock Limits of spiked random matrices {I}.
\newblock {\em Probab.~Theory Rel.~Fields}, 156:795--825, 2013.

\bibitem{BEY}
Paul Bourgade, L{{\'a}}szl{{\'o}} Erd{\H{o}}s, and Horng-Tzer Yau.
\newblock Universality of general {$\beta$}-ensembles.
\newblock {\em Duke Math. J.}, 163(6):1127--1190, 2014.

\bibitem{DVJ1}
Daryl~J Daley and David Vere-Jones.
\newblock {\em An introduction to the theory of point processes: volume I:
  elementary theory and methods}.
\newblock Springer, 2003.

\bibitem{DVJ2}
Daryl~J Daley and David Vere-Jones.
\newblock {\em An introduction to the theory of point processes: volume II:
  general theory and structure}.
\newblock Springer, 2008.

\bibitem{DLR2021}
David Dereudre, Adrien Hardy, Thomas Lebl{\'e}, and Myl{\`e}ne Ma{\"\i}da.
\newblock {DLR} equations and rigidity for the {S}ine-beta process.
\newblock {\em Communications on Pure and Applied Mathematics}, 74(1):172--222,
  2021.

\bibitem{DM2025}
Laure Dumaz and Martin Malvy.
\newblock Long range correlations of the {S}ine$_\beta$ point process.
\newblock in preparation, 2025.

\bibitem{DE}
Ioana Dumitriu and Alan Edelman.
\newblock Matrix models for beta ensembles.
\newblock {\em J. Math. Phys.}, 43(11):5830--5847, 2002.

\bibitem{ES}
Alan Edelman and Brian~D. Sutton.
\newblock From random matrices to stochastic operators.
\newblock {\em J. Stat. Phys.}, 127(6):1121--1165, 2007.

\bibitem{Forrester_1984}
Peter~J. Forrester.
\newblock Analogues between a quantum many body problem and the log-gas.
\newblock {\em Journal of Physics A: Mathematical and General},
  17(10):2059--2067, July 1984.

\bibitem{Forrester_1992}
Peter~J. Forrester.
\newblock Selberg correlation integrals and the $1/r^2$ quantum many-body
  system.
\newblock {\em Nuclear Physics B}, 388(3):671--699, 1992.

\bibitem{Forrester_1994}
Peter~J. Forrester.
\newblock Addendum to `{S}elberg correlation integrals and the $1/r^2$ quantum
  many body system'.
\newblock {\em Nuclear Physics B}, 416(1):377--385, 1994.

\bibitem{ForBook}
Peter~J. Forrester.
\newblock {\em Log-gases and random matrices}, volume~34 of {\em London
  Mathematical Society Monographs Series}.
\newblock Princeton University Press, Princeton, NJ, 2010.

\bibitem{Forrester_20212}
Peter~J. Forrester.
\newblock Differential identities for the structure function of some random
  matrix ensembles.
\newblock {\em Journal of Statistical Physics}, 183(2), May 2021.

\bibitem{Forrester_2021}
Peter~J. Forrester and Anas~A. Rahman.
\newblock Relations between moments for the {J}acobi and {C}auchy random matrix
  ensembles.
\newblock {\em Journal of Mathematical Physics}, 62(7), July 2021.

\bibitem{FR2012}
Peter~J. Forrester and Eric~M. Rains.
\newblock A {F}uchsian matrix differential equation for {S}elberg correlation
  integrals.
\newblock {\em Communications in Mathematical Physics}, 309(3):771--792, 2012.

\bibitem{Concrete}
Ronald~L. Graham, Donald~E. Knuth, and Oren Patashnik.
\newblock {\em Concrete mathematics}.
\newblock Addison-Wesley Publishing Company, Advanced Book Program, Reading,
  MA, 1989.

\bibitem{Haldane1981}
F.~Duncan~M. Haldane.
\newblock Effective harmonic-fluid approach to low-energy properties of
  one-dimensional quantum fluids.
\newblock {\em Phys. Rev. Lett.}, 47:1840--1843, Dec 1981.

\bibitem{HP2018}
Diane Holcomb and Elliot Paquette.
\newblock The maximum deviation of the {S}ine-beta counting process.
\newblock {\em Electronic Communications in Probability}, 23, 2018.

\bibitem{HV}
Diane Holcomb and Benedek Valk{\'o}.
\newblock Large deviations for the {Sine$_\beta$} and {Sch$_\tau$} processes.
\newblock {\em Probability Theory and Related Fields}, 163(1-2):339--378, 2015.

\bibitem{HV2}
Diane Holcomb and Benedek Valk{\'o}.
\newblock Overcrowding asymptotics for the {S}ine$_\beta$ process.
\newblock {\em Annales de l’Institut Henri Poincar{\'e}-Probabilit{\'e}s et
  Statistiques}, 53(3):1181--1195, 2017.

\bibitem{Johansson1998}
Kurt Johansson.
\newblock On fluctuations of eigenvalues of random {H}ermitian matrices.
\newblock {\em Duke Math. J.}, 91(1):151--204, 1998.

\bibitem{KS}
Rowan Killip and Mihai Stoiciu.
\newblock Eigenvalue statistics for {CMV} matrices: from {P}oisson to clock via
  random matrix ensembles.
\newblock {\em Duke Math. J.}, 146(3):361--399, 2009.

\bibitem{KVV}
Eugene Kritchevski, Benedek Valk{{\'o}}, and B{{\'a}}lint Vir{{\'a}}g.
\newblock The scaling limit of the critical one-dimensional random
  {S}chr{\"o}dinger operator.
\newblock {\em Comm. Math. Phys.}, 314(3):775--806, 2012.

\bibitem{Lambert2021}
Gaultier Lambert.
\newblock Mesoscopic central limit theorem for the circular {$\beta$}-ensembles
  and applications.
\newblock {\em Electron. J. Probab.}, 26:Paper No. 7, 33, 2021.

\bibitem{Leble2021}
Thomas Lebl{\'e}.
\newblock {CLT} for fluctuations of linear statistics in the {S}ine-beta
  process.
\newblock {\em International Mathematics Research Notices}, 2021(8):5676--5756,
  2021.

\bibitem{LV_HP}
Yun Li and Benedek Valk{\'o}.
\newblock Operator level limit of the circular {Jacobi} $\beta$-ensemble.
\newblock {\em Random Matrices: Theory and Applications}, 11(04), July 2022.

\bibitem{mehta}
Madan~Lal Mehta.
\newblock {\em Random matrices}, volume 142 of {\em Pure and Applied
  Mathematics (Amsterdam)}.
\newblock Elsevier/Academic Press, Amsterdam, third edition, 2004.

\bibitem{Nakano}
Fumihiko Nakano.
\newblock Level statistics for one-dimensional {S}chr{\"o}dinger operators and
  {G}aussian beta ensemble.
\newblock {\em Journal of Statistical Physics}, 156(1):66--93, 2014.

\bibitem{okounkov1997}
Andrei Okounkov.
\newblock On n-point correlations in the log-gas at rational temperature.
\newblock \url{https://arxiv.org/abs/hep-th/9702001}, 1997.

\bibitem{Protter}
Philip~E. Protter.
\newblock {\em Stochastic integration and differential equations}, volume~21 of
  {\em Stochastic Modelling and Applied Probability}.
\newblock Springer-Verlag, Berlin, 2005.
\newblock Second edition. Version 2.1, Corrected third printing.

\bibitem{RR}
Jos{{\'e}}~A. Ram{\'{\i}}rez and Brian Rider.
\newblock Diffusion at the random matrix hard edge.
\newblock {\em Comm. Math. Phys.}, 288(3):887--906, 2009.

\bibitem{RRV}
Jos{{\'e}}~A. Ram{\'{\i}}rez, Brian Rider, and B{{\'a}}lint Vir{{\'a}}g.
\newblock Beta ensembles, stochastic {A}iry spectrum, and a diffusion.
\newblock {\em J. Amer. Math. Soc.}, 24(4):919--944, 2011.

\bibitem{RRZ}
Jos\'e{}~A. Ram\'irez, Brian Rider, and Ofer Zeitouni.
\newblock Hard edge tail asymptotics.
\newblock {\em Electron. Commun. Probab.}, 16:741--752, 2011.

\bibitem{TW1994}
Craig~A Tracy and Harold Widom.
\newblock Level-spacing distributions and the {A}iry kernel.
\newblock {\em Communications in Mathematical Physics}, 159(1):151--174, 1994.

\bibitem{TW1996}
Craig~A Tracy and Harold Widom.
\newblock On orthogonal and symplectic matrix ensembles.
\newblock {\em Communications in Mathematical Physics}, 177(3):727--754, 1996.

\bibitem{BVBV}
Benedek Valk\'o and B\'alint Vir\'ag.
\newblock Continuum limits of random matrices and the {B}rownian carousel.
\newblock {\em Inventiones Math.}, 177:463--508, 2009.

\bibitem{BVBV2}
Benedek Valk{{\'o}} and B{{\'a}}lint Vir{{\'a}}g.
\newblock Large gaps between random eigenvalues.
\newblock {\em Ann. Probab.}, 38(3):1263--1279, 2010.

\bibitem{BVBV_sbo}
Benedek Valk{\'o} and B{\'a}lint Vir{\'a}g.
\newblock The {Sine$_\beta$} operator.
\newblock {\em Inventiones mathematicae}, 209(1):275--327, 2017.

\bibitem{BVBV_Palm}
Benedek Valk{\'o} and B{\'a}lint Vir{\'a}g.
\newblock Palm measures for {D}irac operators and the {S}ine$_\beta$ process.
\newblock {\em Stochastic Processes and their Applications}, 163:106--135,
  2023.

\end{thebibliography}



\end{document}
